% =====================================================================
% Cahier de travail autonome -- module de mathématiques (niveau licence 1, filières pharma/ingénieur)
% pdfLaTeX + mathpazo, deux passes. Numérotation entièrement automatique.
% Tout le contenu (cours, exemples, exercices, corrigés) est original.
% =====================================================================
\documentclass[10.5pt,a4paper]{article}
\usepackage[utf8]{inputenc}
\usepackage[T1]{fontenc}
\usepackage{mathpazo}
\IfFileExists{french.ldf}{\usepackage[french]{babel}}{\usepackage[english]{babel}}
\usepackage[margin=2cm]{geometry}
\usepackage{amsmath,amssymb}
\usepackage[table]{xcolor}
\usepackage{enumitem}
\usepackage{array,longtable,booktabs}
\usepackage{needspace}
\usepackage[framemethod=default]{mdframed}
\usepackage{fancyhdr}
\usepackage{tikz}

\newif\ifcorr
\corrtrue % \corrfalse : énoncés seuls (sans les corrigés détaillés)

\definecolor{marine}{HTML}{1F3A5F}
\definecolor{acier}{HTML}{4A6FA5}
\definecolor{ardoise}{HTML}{5B6770}
\definecolor{fondpale}{HTML}{F6F8FA}
\definecolor{fondcas}{HTML}{F3F1EE}

\setlist[enumerate,1]{label=\arabic*.,leftmargin=*,itemsep=1pt,topsep=2pt}
\setlist[enumerate,2]{label=\alph*.,leftmargin=*,itemsep=1pt,topsep=1pt}
\setlist[itemize]{leftmargin=*,itemsep=1pt,topsep=2pt}
\sloppy\emergencystretch=3em\setlength{\parindent}{0pt}\setlength{\parskip}{3pt}
\renewcommand{\arraystretch}{1.15}\setlength{\tabcolsep}{5pt}

% --- Compteurs (jamais de numéro saisi à la main) --------------------
\newcounter{defn}[section]\newcounter{prop}[section]\newcounter{demo}[section]
\newcounter{exemple}[section]\newcounter{exo}[section]

\mdfdefinestyle{filet}{linecolor=acier,linewidth=2.5pt,topline=false,bottomline=false,rightline=false,
  leftmargin=0pt,rightmargin=0pt,innerleftmargin=9pt,innerrightmargin=6pt,innertopmargin=3pt,innerbottommargin=3pt,backgroundcolor=fondpale}
\newmdenv[style=filet,linecolor=marine]{boussole}
\newmdenv[style=filet]{methode}
\newmdenv[style=filet,linecolor=ardoise,backgroundcolor=fondcas]{piege}
\newmdenv[style=filet,linecolor=ardoise,backgroundcolor=fondcas]{cadrecas}

\newenvironment{definition}[1][]{\refstepcounter{defn}\par\smallskip\noindent
 \textbf{\textcolor{marine}{Définition \thesection.\thedefn}}\if\relax\detokenize{#1}\relax\else\ (#1)\fi\,\textbf{.}\ }{\par\smallskip}
\newenvironment{propriete}[1][]{\refstepcounter{prop}\par\smallskip\noindent
 \textbf{\textcolor{marine}{Propriété \thesection.\theprop}}\if\relax\detokenize{#1}\relax\else\ (#1)\fi\,\textbf{.}\ \itshape}{\par\smallskip}
\newenvironment{demonstration}[1][]{\par\smallskip\noindent\textbf{\textcolor{ardoise}{Démonstration}}%
 \if\relax\detokenize{#1}\relax\else\ (#1)\fi\,\textbf{.}\ }{\hfill$\square$\par\smallskip}
% Exemple travaillé et cas limite partagent le compteur (ordre de lecture = ordre de difficulté)
\newenvironment{exemple}[1][]{\refstepcounter{exemple}\par\medskip\noindent
 \textbf{\textcolor{marine}{Exemple \thesection.\theexemple}}\if\relax\detokenize{#1}\relax\else\ -- #1\fi\,\textbf{.}\ }{\par\smallskip}
\newenvironment{cas}[1][]{\refstepcounter{exemple}\par\medskip\begin{cadrecas}\noindent
 \textbf{\textcolor{ardoise}{Cas limite \thesection.\theexemple}}\if\relax\detokenize{#1}\relax\else\ -- #1\fi\,\textbf{.}\ }{\end{cadrecas}}
\newenvironment{exo}[1][A]{\refstepcounter{exo}\par\medskip\noindent
 \textbf{\textcolor{marine}{Exercice \thesection.\theexo}}\ \textsf{\small[niveau #1]}\,\textbf{.}\ }{\par}
\newenvironment{corrige}[1]{\par\smallskip\noindent\textbf{\textcolor{acier}{Corrigé de l'exercice #1.}}\ }{\par\smallskip}

\providecommand{\og}{\guillemotleft\,}\providecommand{\fg}{\,\guillemotright}
\newcommand{\R}{\mathbb{R}}\newcommand{\C}{\mathbb{C}}\newcommand{\N}{\mathbb{N}}
\newcommand{\dd}{\mathrm{d}}\newcommand{\e}{\mathrm{e}}
\newcommand{\bareme}[1]{\textsf{\small[#1\,pts]}}
\newcommand{\verif}{\textsf{\textcolor{ardoise}{(convention)}}}
\newcommand{\horsprog}{\textsf{\textcolor{ardoise}{[HORS PROGRAMME]}}}
\newcommand{\titrecourant}{Cahier de travail}
\pagestyle{fancy}\fancyhf{}\renewcommand{\headrulewidth}{0.4pt}
\fancyhead[L]{\small\textcolor{ardoise}{Cahier de travail autonome}}\fancyhead[R]{\small\textcolor{ardoise}{\titrecourant}}
\fancyfoot[C]{\small\thepage}
\usepackage[hidelinks]{hyperref}

\begin{document}
\begin{center}
 {\LARGE\bfseries\textcolor{marine}{Mathématiques : cahier de travail autonome}}\\[3pt]
 {\large\textcolor{ardoise}{Sept semaines : cours, exemples travaillés, cas limites, exercices corrigés}}
\end{center}

\section*{Mode d'emploi}\addcontentsline{toc}{section}{Mode d'emploi}
\begin{boussole}
\textbf{Ce que contient ce cahier.} Sept chapitres indépendants de toute autre ressource. Chacun suit le même parcours, en six temps :
\begin{enumerate}
 \item \textbf{Boussole} : l'objectif de la semaine en une phrase, les compétences vérifiables, le temps estimé.
 \item \textbf{Test d'entrée} : cinq questions pour savoir si les prérequis sont solides (réponses en fin de chapitre).
 \item \textbf{Cours} : définitions, propriétés, démonstrations, méthodes encadrées et pièges.
 \item \textbf{Exemples et cas limites} : par difficulté croissante ; les cas limites (cadre grisé) montrent où une règle cesse de s'appliquer.
 \item \textbf{Exercices} : niveaux A (application directe), B (entraînement), C (approfondissement), tous corrigés en détail.
 \item \textbf{Contrôle de fin de semaine} : trente minutes, vingt points, barème et corrigé.
\end{enumerate}
\textbf{Conseil de travail.} Cherchez chaque exercice sur papier avant d'ouvrir le corrigé ; vérifiez ensuite votre résultat par un test numérique (point particulier, cas limite, ordre de grandeur), comme le font les corrigés.
\end{boussole}
\textbf{Conventions.} Réels $\R$, complexes $\C$, virgule décimale. Une démonstration se termine par $\square$. Les mentions \verif\ signalent un point dont la convention peut varier selon votre cours. Tous les exercices et les valeurs numériques sont originaux ; chaque résultat chiffré a été recalculé par un logiciel de calcul formel avant d'être imprimé.

\renewcommand{\contentsname}{Sommaire}\tableofcontents

\clearpage
\section{Trigonométrie, fonctions usuelles, nombres complexes}
\renewcommand{\titrecourant}{Semaine 1}

\begin{boussole}
\textbf{Objectif de la semaine.} Savoir simplifier une expression trigonométrique ou hyperbolique, convertir un complexe entre forme algébrique et forme exponentielle, résoudre $z^n=a$ et linéariser un produit de $\cos$ et $\sin$, en contrôlant chaque résultat par un test.\\
\textbf{À la fin, je sais :}
\begin{enumerate}
 \item lire sur le cercle valeurs, signes et symétries, sans calculatrice ;
 \item démontrer et employer $\cos(a+b)$, $\cos^2+\sin^2=1$, $\cosh^2-\sinh^2=1$ ;
 \item donner module et argument d'un complexe, quadrant compris ;
 \item résoudre $z^n=a$ et placer les solutions sur le plan complexe ;
 \item compter les zéros terminaux de $n!$ et simplifier des quotients de factorielles.
\end{enumerate}
\textbf{Temps estimé} : 5 à 6 heures (cours 1\,h\,30, exemples 1\,h, exercices 3\,h, contrôle 30\,min).
\end{boussole}

\subsection{Test d'entrée (5 minutes)}
\begin{enumerate}
 \item Convertir $\tfrac{3\pi}{4}$ radians en degrés.
 \item Développer $(a+b)^3$.
 \item Que valent $i^2$ et $i^3$ ?
 \item Simplifier $\ln(\e^{5})$ et $\e^{\ln 7}$.
 \item Donner $\sin\tfrac\pi6$ et $\cos\tfrac\pi3$.
\end{enumerate}
\textit{Réponses : $135^\circ$ ; $a^3+3a^2b+3ab^2+b^3$ ; $-1$ et $-i$ ; $5$ et $7$ ; $\tfrac12$ et $\tfrac12$. Deux erreurs ou plus : refaire le rappel correspondant avant le cours.}

\subsection{Cours}
\subsubsection*{Cercle trigonométrique}
\begin{definition}[Cosinus, sinus, tangente]
Pour un réel $\theta$ (en radians), le point du cercle de centre $O$ et de rayon $1$ repéré par l'angle $\theta$ a pour coordonnées $(\cos\theta,\sin\theta)$. On pose $\tan\theta=\dfrac{\sin\theta}{\cos\theta}$, \emph{définie seulement si $\cos\theta\ne0$}, c'est-à-dire $\theta\not\equiv\frac\pi2\ [\pi]$.
\end{definition}
\begin{propriete}[Relation fondamentale et symétries]
$\cos^2\theta+\sin^2\theta=1$ ; $\cos(-\theta)=\cos\theta$, $\sin(-\theta)=-\sin\theta$ ; $\cos(\pi-\theta)=-\cos\theta$, $\sin(\pi-\theta)=\sin\theta$ ; $\cos(\theta+2\pi)=\cos\theta$ ; $\cos(\tfrac\pi2-\theta)=\sin\theta$.
\end{propriete}
\begin{center}\begin{tabular}{c|ccccc}
$\theta$ & $0$ & $\frac\pi6$ & $\frac\pi4$ & $\frac\pi3$ & $\frac\pi2$\\ \hline
$\cos\theta$ & $1$ & $\frac{\sqrt3}2$ & $\frac{\sqrt2}2$ & $\frac12$ & $0$\\
$\sin\theta$ & $0$ & $\frac12$ & $\frac{\sqrt2}2$ & $\frac{\sqrt3}2$ & $1$\\
$\tan\theta$ & $0$ & $\frac{\sqrt3}3$ & $1$ & $\sqrt3$ & \textit{non définie}
\end{tabular}\end{center}
\begin{definition}[Exponentielle complexe]
Pour $\theta\in\R$, $\e^{i\theta}=\cos\theta+i\sin\theta$. On admet que $\e^{i(a+b)}=\e^{ia}\e^{ib}$.
\end{definition}
\begin{propriete}[Formules d'addition et de duplication]
$\cos(a+b)=\cos a\cos b-\sin a\sin b$ ; $\sin(a+b)=\sin a\cos b+\cos a\sin b$ ; $\cos2a=\cos^2a-\sin^2a=2\cos^2a-1=1-2\sin^2a$ ; $\sin2a=2\sin a\cos a$.
\end{propriete}
\begin{demonstration}
$\e^{i(a+b)}=(\cos a+i\sin a)(\cos b+i\sin b)=(\cos a\cos b-\sin a\sin b)+i(\sin a\cos b+\cos a\sin b)$. Le premier membre vaut $\cos(a+b)+i\sin(a+b)$ ; on identifie parties réelle et imaginaire (les quatre nombres $\cos,\sin$ sont réels). On pose ensuite $b=a$, puis on remplace $\sin^2=1-\cos^2$ ou $\cos^2=1-\sin^2$.
\end{demonstration}

\subsubsection*{Logarithme, exponentielle, fonctions hyperboliques}
\begin{propriete}[Logarithme]
Pour $a,b>0$ : $\ln(ab)=\ln a+\ln b$, $\ln\frac ab=\ln a-\ln b$, $\ln(a^p)=p\ln a$. En revanche $\ln(a+b)$ ne se simplifie pas. Le domaine de $\ln$ est $]0;+\infty[$.
\end{propriete}
\begin{definition}[Cosinus et sinus hyperboliques]
$\cosh x=\dfrac{\e^x+\e^{-x}}2$ et $\sinh x=\dfrac{\e^x-\e^{-x}}2$, pour tout réel $x$.
\end{definition}
\begin{propriete}
$\cosh^2x-\sinh^2x=1$ ; $\cosh$ est paire, $\sinh$ est impaire ; $\cosh x\ge1$ ; $\sinh$ est strictement croissante et bijective de $\R$ sur $\R$.
\end{propriete}
\begin{demonstration}
$\cosh^2x-\sinh^2x=\frac14\big[(\e^{2x}+2+\e^{-2x})-(\e^{2x}-2+\e^{-2x})\big]=1$. Pour $\cosh x\ge1$ : $\cosh x-1=\frac{(\e^{x/2}-\e^{-x/2})^2}2\ge0$.
\end{demonstration}
\begin{methode}
\textbf{Méthode (équation hyperbolique).} Poser $t=\e^x>0$, multiplier par $t$, résoudre l'équation du second degré en $t$, \emph{rejeter} toute racine $t\le0$, puis revenir à $x=\ln t$.
\end{methode}

\subsubsection*{Factorielles}
\begin{definition}
$0!=1$ et $n!=n\,(n-1)!=1\cdot2\cdots n$ pour $n\in\N^*$. Ainsi $\dfrac{(n+k)!}{n!}=(n+1)(n+2)\cdots(n+k)$.
\end{definition}
\begin{propriete}[Zéros terminaux]
Le nombre de zéros par lesquels se termine l'écriture décimale de $n!$ est $\displaystyle\sum_{k\ge1}\Big\lfloor\frac n{5^k}\Big\rfloor$.
\end{propriete}
\begin{demonstration}
Un zéro final provient d'un facteur $10=2\cdot5$ ; dans $n!$ les facteurs $2$ sont plus nombreux que les facteurs $5$, donc seul le nombre de facteurs $5$ compte. Il y a $\lfloor n/5\rfloor$ multiples de $5$ entre $1$ et $n$ ; $\lfloor n/25\rfloor$ d'entre eux apportent un deuxième facteur $5$, $\lfloor n/125\rfloor$ un troisième, etc.
\end{demonstration}

\subsubsection*{Nombres complexes}
\begin{definition}[Module, argument]
Pour $z=a+ib\neq0$, $|z|=\sqrt{a^2+b^2}$ et un argument $\theta$ est un réel tel que $\cos\theta=\frac a{|z|}$ et $\sin\theta=\frac b{|z|}$ ; il est défini modulo $2\pi$. On écrit $z=|z|\,\e^{i\theta}$. Le nombre $0$ n'a pas d'argument.
\end{definition}
\begin{propriete}
$|z_1z_2|=|z_1||z_2|$ et $\arg(z_1z_2)=\arg z_1+\arg z_2\ [2\pi]$ ; $\arg\bar z=-\arg z$ ; formule de Moivre : $(\e^{i\theta})^n=\e^{in\theta}$ ; formules d'Euler : $\cos\theta=\frac{\e^{i\theta}+\e^{-i\theta}}2$, $\sin\theta=\frac{\e^{i\theta}-\e^{-i\theta}}{2i}$.
\end{propriete}
\begin{propriete}[Racines $n$-ièmes]
Soit $a=r\e^{i\theta}\ne0$ et $n\ge1$. L'équation $z^n=a$ possède exactement $n$ solutions : $z_k=r^{1/n}\,\e^{i(\theta+2k\pi)/n}$, $k=0,\dots,n-1$. Elles sont les sommets d'un polygone régulier de centre $O$. Pour $a=0$, l'unique solution est $z=0$.
\end{propriete}
\begin{demonstration}
Si $z=\rho\e^{i\varphi}$, alors $z^n=a$ équivaut à $\rho^n=r$ et $n\varphi\equiv\theta\ [2\pi]$, c'est-à-dire $\varphi=\frac{\theta+2k\pi}n$ ; $k$ et $k+n$ donnent la même solution, d'où $n$ solutions distinctes.
\end{demonstration}
\begin{propriete}[Somme des racines de l'unité]
Pour $n\ge2$ et $\omega=\e^{2i\pi/n}$ : $\sum_{k=0}^{n-1}\omega^k=0$. Pour $n=1$, la somme vaut $1$.
\end{propriete}
\begin{demonstration}
$\omega\ne1$ si $n\ge2$, donc $\sum_{k=0}^{n-1}\omega^k=\dfrac{1-\omega^n}{1-\omega}=0$ car $\omega^n=1$.
\end{demonstration}
\begin{methode}
\textbf{M1 -- Argument.} Calculer $|z|$, écrire $\cos\theta=a/|z|$ et $\sin\theta=b/|z|$, reconnaître l'angle dans le tableau des valeurs remarquables, \emph{contrôler le quadrant}. \textbf{M2 -- Linéarisation.} Remplacer $\cos x,\sin x$ par les formules d'Euler, développer par le binôme, regrouper $\e^{ikx}+\e^{-ikx}=2\cos kx$. \textbf{M3 -- $z^n=a$.} Écrire $a$ sous forme exponentielle, appliquer la propriété des racines, vérifier une racine en l'élevant à la puissance $n$.
\end{methode}
\begin{piege}
\textbf{Pièges fréquents.} (i) $\arg(a+ib)=\arctan\frac ba$ est faux dès que $a<0$ : l'arc tangente est dans $]-\frac\pi2;\frac\pi2[$. (ii) $\ln(a+b)\ne\ln a+\ln b$. (iii) $\sqrt{x^2}=|x|$. (iv) $\cosh$ n'est pas périodique. (v) $\tan$ n'est pas définie en $\frac\pi2+k\pi$.
\end{piege}

\subsection{Exemples travaillés et cas limites}
\begin{exemple}[valeurs remarquables]
Calculer $\cos\frac{7\pi}6$, $\sin(-\frac\pi3)$, $\tan\frac{3\pi}4$.\\
\emph{Solution.} $\frac{7\pi}6=\pi+\frac\pi6$ : le point est opposé à celui d'angle $\frac\pi6$, donc $\cos\frac{7\pi}6=-\frac{\sqrt3}2$. Par imparité, $\sin(-\frac\pi3)=-\frac{\sqrt3}2$. Enfin $\tan\frac{3\pi}4=\frac{\sin(\pi/4)}{-\cos(\pi/4)}=-1$.
\end{exemple}
\begin{exemple}[identité sur un domaine]
Exprimer $\tan^2x$ en fonction de $\cos x$.\\
\emph{Solution.} $\tan^2x=\dfrac{\sin^2x}{\cos^2x}=\dfrac{1-\cos^2x}{\cos^2x}=\dfrac1{\cos^2x}-1$.
\end{exemple}
\begin{cas}[domaine de validité]
L'identité précédente n'a de sens que si $\cos x\ne0$, soit $x\ne\frac\pi2+k\pi$. Tester : $x=\frac\pi4$ donne $\tan^2=1$ et $\frac1{1/2}-1=1$. Pour $x=\frac\pi2$ les deux membres sont non définis : on ne doit pas \og vérifier \fg{} en ce point.
\end{cas}
\begin{exemple}[échelle contre un mur]
Une échelle de $5\,\mathrm m$ a son pied à $1{,}4\,\mathrm m$ du mur. Quelle hauteur atteint-elle, et quel angle fait-elle avec le sol ?\\
\emph{Solution.} Théorème de Pythagore : $h=\sqrt{5^2-1{,}4^2}=\sqrt{23{,}04}=4{,}8\,\mathrm m$. L'angle $\alpha$ avec le sol vérifie $\cos\alpha=\frac{1{,}4}5=0{,}28$, donc $\alpha=\arccos0{,}28\approx73{,}74^\circ$. \emph{Contrôle} : $\sin\alpha=\frac{4{,}8}5=0{,}96$ et $0{,}28^2+0{,}96^2=1$.
\end{exemple}
\begin{cas}[bornes physiques]
Le pied $d$ est dans $[0\,;5]$. Si $d=0$ l'échelle est verticale ($h=5$, $\alpha=90^\circ$) ; si $d=5$ elle est couchée ($h=0$). Si le sommet descend de $0{,}4\,\mathrm m$ ($h=4{,}4$), le pied recule à $\sqrt{25-19{,}36}=\sqrt{5{,}64}\approx2{,}37\,\mathrm m$ : un petit glissement vertical près du sommet vertical déplace beaucoup le pied.
\end{cas}
\begin{exemple}[argument d'un complexe]
Forme exponentielle de $z=-1+i\sqrt3$.\\
\emph{Solution.} $|z|=\sqrt{1+3}=2$, $\cos\theta=-\frac12$, $\sin\theta=\frac{\sqrt3}2$ donc $\theta=\frac{2\pi}3$ et $z=2\e^{2i\pi/3}$. La formule naïve $\arctan(\sqrt3/(-1))=-\frac\pi3$ donne un angle faux (décalé de $\pi$).
\end{exemple}
\begin{cas}[arguments extrêmes]
$z=-i$ : $|z|=1$, $\cos\theta=0$, $\sin\theta=-1$, donc $\theta=-\frac\pi2$ (et non \og $\arctan(-1/0)$ \fg{}). $z=-4$ : $\theta=\pi$. $z=0$ : aucun argument. Un réel positif a pour argument $0$ modulo $2\pi$.
\end{cas}
\begin{exemple}[racines cubiques]
Résoudre $z^3=8i$.\\
\emph{Solution.} $8i=8\e^{i\pi/2}$, donc $z_k=2\e^{i(\pi/6+2k\pi/3)}$ : $z_0=2\e^{i\pi/6}=\sqrt3+i$, $z_1=2\e^{i5\pi/6}=-\sqrt3+i$, $z_2=2\e^{i3\pi/2}=-2i$. \emph{Contrôle} : $(-2i)^3=-8i^3=8i$. $\checkmark$
\end{exemple}
\begin{cas}[second membre nul ou réel]
$z^4=0$ n'a que $z=0$. $z^3=8$ a une racine réelle ($2$) et deux complexes conjuguées $-1\pm i\sqrt3$. $z^2=-9$ donne $z=\pm3i$ : la \og racine carrée d'un nombre négatif \fg{} n'est pas un réel.
\end{cas}
\begin{exemple}[linéarisation et intégrale]
Linéariser $\cos^3x$ puis calculer $\int_0^{\pi/2}\cos^3x\,\dd x$.\\
\emph{Solution.} $\cos^3x=\frac18(\e^{ix}+\e^{-ix})^3=\frac18(\e^{3ix}+3\e^{ix}+3\e^{-ix}+\e^{-3ix})=\frac{\cos3x+3\cos x}4$. Alors $\int_0^{\pi/2}\cos^3x\,\dd x=\Big[\frac{\sin3x}{12}+\frac{3\sin x}4\Big]_0^{\pi/2}=-\frac1{12}+\frac34=\frac23$.
\end{exemple}
\begin{exemple}[équation hyperbolique]
Résoudre $\cosh x=2$.\\
\emph{Solution.} Avec $t=\e^x$ : $t+\frac1t=4$, soit $t^2-4t+1=0$, $t=2\pm\sqrt3$, tous deux positifs. Donc $x=\pm\ln(2+\sqrt3)$ (car $(2+\sqrt3)(2-\sqrt3)=1$).
\end{exemple}
\begin{cas}[nombre de solutions]
$\cosh x=1$ n'a que $x=0$ ; $\cosh x=\frac12$ n'a aucune solution (car $\cosh\ge1$) ; en revanche $\sinh x=c$ a toujours exactement une solution $x=\ln(c+\sqrt{c^2+1})$.
\end{cas}
\begin{exemple}[zéros de $n!$]
Combien de zéros terminent $50!$ ?\quad $\lfloor\frac{50}5\rfloor+\lfloor\frac{50}{25}\rfloor=10+2=12$.
\end{exemple}
\begin{cas}[seuils]
$n<5$ : aucun zéro. $n=25$ : $5+1=6$ (le facteur $25$ compte deux fois). $n=124$ : $24+4=28$ ; $n=125$ : $25+5+1=31$ : le nombre de zéros \emph{saute} de trois en passant de $124$ à $125$.
\end{cas}

\subsection{Exercices}
\begin{exo}[A]\label{c1e1}
Calculer sans calculatrice $\cos\frac{5\pi}6$, $\sin\frac{7\pi}4$ et $\tan(-\frac\pi6)$.
\end{exo}
\begin{exo}[A]\label{c1e2}
Simplifier $\dfrac{(n+2)!}{(n-1)!}$ pour $n\ge1$, puis $\dfrac{(2n)!}{(2n-2)!}$ pour $n\ge1$.
\end{exo}
\begin{exo}[A]\label{c1e3}
Une échelle de $6{,}5\,\mathrm m$ a son pied à $2{,}5\,\mathrm m$ du mur. (a) Quelle hauteur atteint-elle ? (b) Quel angle fait-elle avec le sol ? (c) À quelle distance placer le pied pour atteindre exactement $6{,}4\,\mathrm m$ ?
\end{exo}
\begin{exo}[B]\label{c1e4}
Démontrer que $\cos3x=4\cos^3x-3\cos x$. En déduire que $\cos\frac\pi9$ est solution de $8u^3-6u-1=0$.
\end{exo}
\begin{exo}[B]\label{c1e5}
Résoudre dans $[0;2\pi[$ l'équation $\sin x=\cos2x$.
\end{exo}
\begin{exo}[B]\label{c1e6}
Mettre sous forme algébrique $z=\dfrac{(1+i)^8}{(1-i\sqrt3)^3}$.
\end{exo}
\begin{exo}[B]\label{c1e7}
Résoudre dans $\C$ : (a) $z^2+2z+5=0$ ; (b) $z^4=-16$.
\end{exo}
\begin{exo}[C]\label{c1e8}
Résoudre dans $\R$ : $3\cosh x-\sinh x=3$.
\end{exo}
\begin{exo}[C]\label{c1e9}
Soit $n\ge2$ et $\omega_k=\e^{2ik\pi/n}$. (a) Montrer que $\sum_{k=0}^{n-1}\cos\frac{2k\pi}n=0$. (b) Montrer que le produit des $\omega_k$ vaut $(-1)^{n+1}$. (c) Que se passe-t-il pour $n=1$ ?
\end{exo}
\begin{exo}[C]\label{c1e10}
Linéariser $\sin^2x\cos^2x$ puis calculer $\int_0^\pi\sin^2x\cos^2x\,\dd x$.
\end{exo}

\ifcorr
\subsection{Corrigés détaillés}
\begin{corrige}{\ref{c1e1}}
$\frac{5\pi}6=\pi-\frac\pi6$ donc $\cos\frac{5\pi}6=-\cos\frac\pi6=-\frac{\sqrt3}2$. $\frac{7\pi}4=2\pi-\frac\pi4$ donc $\sin\frac{7\pi}4=-\sin\frac\pi4=-\frac{\sqrt2}2$. $\tan(-\frac\pi6)=-\tan\frac\pi6=-\frac{\sqrt3}3$.
\end{corrige}
\begin{corrige}{\ref{c1e2}}
$\frac{(n+2)!}{(n-1)!}=(n+2)(n+1)n$ (trois facteurs, de $n$ à $n+2$). Pour $\frac{(2n)!}{(2n-2)!}$, il reste les facteurs $2n$ et $2n-1$ : $2n(2n-1)$. \emph{Test} : $n=1$ donne $\frac{3!}{0!}=6=3\cdot2\cdot1$ et $\frac{2!}{0!}=2=2\cdot1$. Le cas $n=1$ montre pourquoi $0!=1$ est nécessaire.
\end{corrige}
\begin{corrige}{\ref{c1e3}}
(a) $h=\sqrt{6{,}5^2-2{,}5^2}=\sqrt{42{,}25-6{,}25}=\sqrt{36}=6\,\mathrm m$. (b) $\cos\alpha=\frac{2{,}5}{6{,}5}=\frac5{13}$, donc $\alpha\approx67{,}38^\circ$ ; contrôle : $\sin\alpha=\frac6{6{,}5}=\frac{12}{13}$ et $\frac{25+144}{169}=1$. (c) $d=\sqrt{42{,}25-6{,}4^2}=\sqrt{42{,}25-40{,}96}=\sqrt{1{,}29}\approx1{,}14\,\mathrm m$. Cas limite : $6{,}5$ m serait la hauteur maximale (échelle verticale) ; viser $6{,}6\,\mathrm m$ est impossible, ce que révèle un radicande négatif.
\end{corrige}
\begin{corrige}{\ref{c1e4}}
$\cos3x=\cos(2x+x)=\cos2x\cos x-\sin2x\sin x=(2\cos^2x-1)\cos x-2\sin^2x\cos x=(2\cos^2x-1)\cos x-2(1-\cos^2x)\cos x=4\cos^3x-3\cos x$. Avec $x=\frac\pi9$ : $\cos\frac\pi3=\frac12=4u^3-3u$ où $u=\cos\frac\pi9$, soit $8u^3-6u-1=0$. \emph{Test} : $u\approx0{,}9397$ donne $8\cdot0{,}8298-5{,}638-1\approx0$.
\end{corrige}
\begin{corrige}{\ref{c1e5}}
$\cos2x=1-2\sin^2x$, donc l'équation devient $2s^2+s-1=0$ avec $s=\sin x$, soit $(2s-1)(s+1)=0$, $s=\frac12$ ou $s=-1$. Dans $[0;2\pi[$ : $s=\frac12$ donne $x=\frac\pi6$ ou $\frac{5\pi}6$ ; $s=-1$ donne $x=\frac{3\pi}2$. Solutions : $\{\frac\pi6,\frac{5\pi}6,\frac{3\pi}2\}$. \emph{Contrôle} en $\frac{3\pi}2$ : $\sin=-1$ et $\cos3\pi=-1$. $\checkmark$
\end{corrige}
\begin{corrige}{\ref{c1e6}}
$1+i=\sqrt2\,\e^{i\pi/4}$, donc $(1+i)^8=(\sqrt2)^8\e^{2i\pi}=16$. $1-i\sqrt3=2\e^{-i\pi/3}$, donc $(1-i\sqrt3)^3=8\e^{-i\pi}=-8$. Donc $z=\frac{16}{-8}=-2$. \emph{Contrôle} : $(1+i)^2=2i$, $(2i)^4=16$ ; $(1-i\sqrt3)^2=-2-2i\sqrt3$, puis $(-2-2i\sqrt3)(1-i\sqrt3)=-2+2i\sqrt3-2i\sqrt3-6=-8$. $\checkmark$
\end{corrige}
\begin{corrige}{\ref{c1e7}}
(a) $\Delta=4-20=-16=(4i)^2$, donc $z=\frac{-2\pm4i}2=-1\pm2i$. (b) $-16=16\e^{i\pi}$, donc $z_k=2\e^{i(\pi/4+k\pi/2)}$, $k=0,\dots,3$ : $\sqrt2(1+i)$, $\sqrt2(-1+i)$, $\sqrt2(-1-i)$, $\sqrt2(1-i)$, c'est-à-dire $\pm\sqrt2\pm i\sqrt2$. Contrôle : $(\sqrt2(1+i))^4=4(1+i)^4=4(2i)^2=-16$. $\checkmark$
\end{corrige}
\begin{corrige}{\ref{c1e8}}
Avec $t=\e^x>0$ : $\frac32(t+\frac1t)-\frac12(t-\frac1t)=3$, soit $t+\frac2t=3$, $t^2-3t+2=0$, $t\in\{1,2\}$, tous deux positifs. Donc $x=0$ ou $x=\ln2$. Contrôle en $\ln2$ : $\cosh=\frac{2+1/2}2=1{,}25$, $\sinh=\frac{2-1/2}2=0{,}75$ ; $3\cdot1{,}25-0{,}75=3$. $\checkmark$
\end{corrige}
\begin{corrige}{\ref{c1e9}}
(a) $\sum\omega_k$ a une partie réelle $\sum\cos\frac{2k\pi}n$ ; or $\sum\omega_k=0$ (propriété du cours), donc cette partie réelle est nulle. (b) $\prod_{k=0}^{n-1}\omega_k=\e^{\frac{2i\pi}n\sum k}=\e^{\frac{2i\pi}n\cdot\frac{n(n-1)}2}=\e^{i\pi(n-1)}=(-1)^{n-1}=(-1)^{n+1}$. Contrôle : $n=2$, racines $1,-1$, produit $-1$ ; $n=3$, produit $1$. (c) Pour $n=1$ la seule racine est $1$ : la somme vaut $1\ne0$ et le produit vaut $1=(-1)^2$ ; la formule du produit reste vraie, celle de la somme exige $n\ge2$.
\end{corrige}
\begin{corrige}{\ref{c1e10}}
$\sin x\cos x=\frac12\sin2x$, donc $\sin^2x\cos^2x=\frac14\sin^22x=\frac14\cdot\frac{1-\cos4x}2=\frac{1-\cos4x}8$. Alors $\int_0^\pi\frac{1-\cos4x}8\dd x=\frac18\big[x-\frac{\sin4x}4\big]_0^\pi=\frac\pi8$. Contrôle numérique : la valeur moyenne de $\sin^2\cos^2$ sur une période est $\frac18$, et $\frac18\cdot\pi=\frac\pi8$. $\checkmark$
\end{corrige}
\else\stepcounter{subsection}\fi

\subsection{Contrôle de fin de semaine (30 minutes, 20 points)}
\begin{enumerate}
 \item \bareme{6} Donner la forme exponentielle de $z=-\sqrt3+i$, puis calculer $z^6$.
 \item \bareme{7} Résoudre dans $[0;2\pi[$ : $2\cos^2x+\cos x-1=0$.
 \item \bareme{7} Résoudre dans $\R$ : $\cosh x=\frac32$ et montrer que les solutions sont $\pm2\ln\varphi$, où $\varphi=\frac{1+\sqrt5}2$.
\end{enumerate}
\ifcorr
\textbf{Corrigé et barème.}
\textbf{1.} $|z|=2$, $\cos\theta=-\frac{\sqrt3}2$, $\sin\theta=\frac12$, donc $\theta=\frac{5\pi}6$ : $z=2\e^{5i\pi/6}$ \bareme{3 : module et argument avec le quadrant}. $z^6=64\e^{5i\pi}=-64$ \bareme{3 : Moivre et résultat}. \textbf{2.} $2c^2+c-1=(2c-1)(c+1)$ avec $c=\cos x$ \bareme{3} ; $c=\frac12$ donne $\frac\pi3,\frac{5\pi}3$ ; $c=-1$ donne $\pi$ \bareme{4 : trois solutions, aucune oubliée}. \textbf{3.} $t=\e^x$ : $t^2-3t+1=0$, $t=\frac{3\pm\sqrt5}2>0$ \bareme{3} ; $x=\ln\frac{3+\sqrt5}2$ ou $\ln\frac{3-\sqrt5}2=-\ln\frac{3+\sqrt5}2$ ; comme $\varphi^2=\frac{6+2\sqrt5}4=\frac{3+\sqrt5}2$, on a $x=\pm2\ln\varphi$ \bareme{4}. Seuil de réussite conseillé : $14/20$.
\fi

\subsection{Corrigé du test d'entrée et grille d'auto-évaluation}
Cochez : $\square$ je lis cercle et symétries ; $\square$ je trouve un argument sans me tromper de quadrant ; $\square$ je résous $z^n=a$ et je vérifie une racine ; $\square$ je linéarise avec Euler ; $\square$ je rejette les racines $t\le0$ d'une équation hyperbolique. Moins de trois cases cochées : reprendre les exemples correspondants avant la Semaine~2.
\clearpage
\section{Proportions, polynômes, fractions rationnelles, limites}
\renewcommand{\titrecourant}{Semaine 2}

\begin{boussole}
\textbf{Objectif de la semaine.} Traduire un énoncé de proportion en équation, résoudre sans erreur les équations du second degré et bicarrées (en sachant combien de solutions attendre), décomposer une fraction rationnelle en éléments simples, et écrire une limite avec la définition $\varepsilon$--$\delta$.\\
\textbf{À la fin, je sais :}
\begin{enumerate}
 \item enchaîner des pourcentages par des facteurs multiplicatifs et inverser une variation ;
 \item interpoler linéairement entre deux mesures, et dire quand ne pas extrapoler ;
 \item discuter le nombre de solutions d'une équation du second degré dépendant d'un paramètre ;
 \item décomposer $\frac{P}{(x-a)(x-b)}$ (avec division euclidienne si le degré l'exige) et traiter un pôle double ;
 \item prouver $\lim_{x\to a}f(x)=\ell$ en exhibant $\delta(\varepsilon)$.
\end{enumerate}
\textbf{Temps estimé} : 5 à 6 heures.
\end{boussole}

\subsection{Test d'entrée (5 minutes)}
\begin{enumerate}
 \item Quel nombre multiplie un prix lors d'une hausse de $20\,\%$ ?
 \item Développer $(x-3)^2$ et factoriser $x^2-9$.
 \item Résoudre $3x-7=2x+5$.
 \item Que vaut le discriminant de $x^2-4x+4$ ?
 \item Réduire au même dénominateur $\frac1x+\frac1{x+1}$.
\end{enumerate}
\textit{Réponses : $1{,}2$ ; $x^2-6x+9$ et $(x-3)(x+3)$ ; $x=12$ ; $0$ ; $\frac{2x+1}{x(x+1)}$.}

\subsection{Cours}
\subsubsection*{Pourcentages et interpolation}
\begin{definition}[Facteur multiplicatif]
Une variation de $t\,\%$ multiplie la valeur par $1+\frac t{100}$ (hausse si $t>0$, baisse si $t<0$). Des variations successives multiplient leurs facteurs ; la variation réciproque est le facteur inverse, \emph{pas} l'opposé du taux.
\end{definition}
\begin{propriete}[Interpolation linéaire]
Si une grandeur $y$ est affine en $x$ et connue en $x_0\ne x_1$ ($y_0$, $y_1$), alors $y(x)=y_0+\dfrac{y_1-y_0}{x_1-x_0}(x-x_0)$. Si $x\in[x_0;x_1]$ on interpole ; sinon on \emph{extrapole}, ce qui suppose que le modèle affine reste valable.
\end{propriete}
\begin{demonstration}
Une fonction affine s'écrit $mx+p$ ; sa pente vaut $m=\frac{y_1-y_0}{x_1-x_0}$ et $y(x)-y(x_0)=m(x-x_0)$.
\end{demonstration}

\subsubsection*{Second degré et équations bicarrées}
\begin{propriete}[Discriminant]
Soit $a\ne0$ et $\Delta=b^2-4ac$. Si $\Delta>0$ : deux racines réelles $\frac{-b\pm\sqrt\Delta}{2a}$. Si $\Delta=0$ : une racine double $-\frac b{2a}$. Si $\Delta<0$ : aucune racine réelle, deux racines complexes conjuguées $\frac{-b\pm i\sqrt{-\Delta}}{2a}$.
\end{propriete}
\begin{demonstration}
Forme canonique : $ax^2+bx+c=a\Big[\big(x+\frac b{2a}\big)^2-\frac\Delta{4a^2}\Big]$. L'équation équivaut à $\big(x+\frac b{2a}\big)^2=\frac\Delta{4a^2}$, d'où les trois cas selon le signe de $\Delta$.
\end{demonstration}
\begin{propriete}[Somme et produit des racines]
Si $x_1,x_2$ sont les racines : $x_1+x_2=-\frac ba$ et $x_1x_2=\frac ca$ ; et $ax^2+bx+c=a(x-x_1)(x-x_2)$.
\end{propriete}
\begin{demonstration}
Développer $a(x-x_1)(x-x_2)=a[x^2-(x_1+x_2)x+x_1x_2]$ et identifier avec $ax^2+bx+c$.
\end{demonstration}
\begin{methode}
\textbf{Bicarrée $ax^4+bx^2+c=0$.} Poser $X=x^2$ ; résoudre $aX^2+bX+c=0$ ; \emph{garder $X\ge0$ seulement} ; chaque $X>0$ donne $x=\pm\sqrt X$, $X=0$ donne $x=0$. Une racine $X<0$ ne donne aucune solution réelle (deux solutions imaginaires pures).
\end{methode}

\subsubsection*{Éléments simples}
\begin{propriete}[Deux pôles simples]
Pour $a\ne b$ et $\deg P\le1$ : $\dfrac{P(x)}{(x-a)(x-b)}=\dfrac A{x-a}+\dfrac B{x-b}$ avec $A=\dfrac{P(a)}{a-b}$ et $B=\dfrac{P(b)}{b-a}$.
\end{propriete}
\begin{demonstration}
En multipliant par $(x-a)(x-b)$ : $P(x)=A(x-b)+B(x-a)$. En $x=a$ : $P(a)=A(a-b)$ ; en $x=b$ : $P(b)=B(b-a)$. Les deux membres sont des polynômes de degré $\le1$ égaux en deux points distincts, donc égaux.
\end{demonstration}
\begin{methode}
(1) Si $\deg P\ge\deg(\text{dénominateur})$, faire d'abord la division euclidienne : le quotient est un polynôme. (2) Pôle double $\frac1{(x-a)^2}$ : on cherche $\frac A{x-a}+\frac B{(x-a)^2}$ ; $B$ s'obtient en évaluant la fraction multipliée par $(x-a)^2$ en $x=a$. (3) \emph{Toujours vérifier} en une valeur numérique de $x$ non pôle.
\end{methode}

\subsubsection*{Limites}
\begin{definition}[Limite en un point]
$\displaystyle\lim_{x\to a}f(x)=\ell$ si $\forall\varepsilon>0,\ \exists\delta>0,\ \forall x,\ 0<|x-a|<\delta\Rightarrow|f(x)-\ell|<\varepsilon$. L'ordre des quantificateurs est essentiel : $\delta$ dépend de $\varepsilon$ (et de $a$), jamais l'inverse. La condition $0<|x-a|$ exclut $x=a$ : la valeur de $f$ en $a$ n'intervient pas.
\end{definition}
\begin{propriete}
Une limite, si elle existe, est unique. Si $f$ a une limite finie en $a$, ses limites à gauche et à droite existent et sont égales ; la réciproque est vraie.
\end{propriete}
\begin{piege}
\textbf{Pièges.} (i) $\varepsilon$ et $\delta$ inversés. (ii) Croire qu'une limite est la valeur $f(a)$. (iii) $\frac00$ n'est pas un nombre : une fraction qui s'annule en haut et en bas se simplifie ou se factorise d'abord. (iv) Multiplier deux pourcentages \og comme on les additionne \fg{}.
\end{piege}

\subsection{Exemples travaillés et cas limites}
\begin{exemple}[pourcentages successifs]
Un prix augmente de $20\,\%$ puis baisse de $20\,\%$. Variation globale ?\\
\emph{Solution.} Facteur $1{,}2\times0{,}8=0{,}96$ : baisse globale de $4\,\%$ (et non $0\,\%$).
\end{exemple}
\begin{cas}[variation réciproque]
Après une baisse de $25\,\%$, un objet coûte $90$. Prix initial : $\frac{90}{0{,}75}=120$ (et non $90\times1{,}25=112{,}5$, qui serait $+25\,\%$ \emph{du prix déjà réduit}). Pour revenir au prix initial il faut une hausse de $\frac1{0{,}75}-1\approx+33{,}3\,\%$.
\end{cas}
\begin{exemple}[interpolation]
Une boisson est à $80^\circ$C à $t=10$ min et à $50^\circ$C à $t=30$ min. Estimer la température à $t=25$.\\
\emph{Solution.} Pente $\frac{50-80}{30-10}=-1{,}5$ ; $T(25)=80-1{,}5\cdot15=57{,}5^\circ$C.
\end{exemple}
\begin{cas}[extrapolation et impossibilité]
Le même modèle donnerait $T(70)=-10^\circ$C, température absurde pour une boisson dans une pièce chaude : le refroidissement réel n'est pas affine. Et si $x_0=x_1$, la formule divise par $0$ : deux mesures au même instant ne déterminent pas de droite.
\end{cas}
\begin{exemple}[discriminant]
Résoudre $2x^2+x-6=0$.\\
\emph{Solution.} $\Delta=1+48=49$, $x=\frac{-1\pm7}4$, soit $x=\frac32$ ou $x=-2$. Contrôle : somme $-\frac12=-\frac ba$ ; produit $-3=\frac ca$. $\checkmark$
\end{exemple}
\begin{cas}[équation à paramètre]
Discuter le nombre de solutions réelles de $(m-1)x^2-2x+1=0$.\\
\emph{Solution.} Si $m=1$, l'équation n'est plus du second degré : $-2x+1=0$, $x=\frac12$ (une solution). Si $m\ne1$ : $\Delta=4-4(m-1)=8-4m$. Donc $m<2$ ($m\ne1$) : deux solutions ; $m=2$ : une solution double $x=1$ ; $m>2$ : aucune. \emph{Contrôle} $m=2$ : $x^2-2x+1=(x-1)^2$. $\checkmark$ Oublier le cas $m=1$ est l'erreur typique.
\end{cas}
\begin{exemple}[bicarrée]
Résoudre $x^4-5x^2+4=0$.\\
\emph{Solution.} $X^2-5X+4=0$ donne $X=1$ ou $X=4$, tous deux positifs : $x\in\{-2,-1,1,2\}$.
\end{exemple}
\begin{cas}[racine $X$ négative]
$x^4+x^2-6=0$ : $X\in\{2,-3\}$. $X=2$ donne $x=\pm\sqrt2$ ; $X=-3$ est rejetée dans $\R$ (dans $\C$ : $x=\pm i\sqrt3$). Enfin $x^4+2x^2+1=(x^2+1)^2$ n'a aucune racine réelle.
\end{cas}
\begin{exemple}[éléments simples]
Décomposer $\dfrac{3x+5}{(x-1)(x+2)}$.\\
\emph{Solution.} $A=\frac{3+5}{1+2}=\frac83$, $B=\frac{-6+5}{-2-1}=\frac13$. \emph{Contrôle} en $x=0$ : membre de gauche $\frac5{-2}=-\frac52$ ; membre de droite $-\frac83+\frac16=-\frac{15}6=-\frac52$. $\checkmark$
\end{exemple}
\begin{cas}[degré trop grand, pôle double]
(a) $\frac{x^2+1}{(x-1)(x+1)}=1+\frac2{(x-1)(x+1)}=1+\frac1{x-1}-\frac1{x+1}$ : sans division euclidienne, la méthode directe donnerait un faux résultat. (b) $\frac x{(x-1)^2}=\frac{(x-1)+1}{(x-1)^2}=\frac1{x-1}+\frac1{(x-1)^2}$. (c) La formule de $A$ et $B$ exige $a\ne b$ : avec $a=b$ elle divise par $0$.
\end{cas}
\begin{exemple}[limite $\varepsilon$--$\delta$ affine]
Montrer $\lim_{x\to2}(3x-1)=5$.\\
\emph{Solution.} $|3x-1-5|=3|x-2|$. Pour $\varepsilon>0$ on prend $\delta=\frac\varepsilon3$ : $|x-2|<\delta\Rightarrow|f(x)-5|<\varepsilon$.
\end{exemple}
\begin{exemple}[limite d'un carré]
Montrer $\lim_{x\to2}x^2=4$.\\
\emph{Solution.} $|x^2-4|=|x-2||x+2|$. Si $|x-2|<1$ alors $x\in]1;3[$ et $|x+2|<5$. Avec $\delta=\min(1,\frac\varepsilon5)$ : $|x^2-4|<5\cdot\frac\varepsilon5=\varepsilon$.
\end{exemple}
\begin{cas}[limite sans valeur, limite inexistante]
$\frac{x^2-1}{x-1}=x+1$ pour $x\ne1$, donc la limite en $1$ vaut $2$ alors que la fonction n'est pas définie en $1$. À l'inverse, $f(x)=\frac{|x|}x$ vaut $-1$ à gauche de $0$ et $1$ à droite : les limites latérales diffèrent, la limite en $0$ n'existe pas.
\end{cas}

\subsection{Exercices}
\begin{exo}[A]\label{c2e1}
Un article à $80$ \texteuro{} augmente de $15\,\%$ puis baisse de $15\,\%$. Prix final et variation globale ? Un autre article coûte $56$ \texteuro{} après une remise de $20\,\%$ : prix initial ?
\end{exo}
\begin{exo}[A]\label{c2e2}
Un capteur indique $12^\circ$C à $t=0$ min et $4^\circ$C à $t=40$ min. En supposant l'évolution affine, estimer la température à $t=25$ min, puis l'instant où elle atteint $0^\circ$C. Commenter.
\end{exo}
\begin{exo}[A]\label{c2e3}
Résoudre $2x^2-5x-3=0$ et factoriser le trinôme.
\end{exo}
\begin{exo}[B]\label{c2e4}
Discuter selon $m$ le nombre de solutions réelles de $x^2-2mx+m+2=0$.
\end{exo}
\begin{exo}[B]\label{c2e5}
Résoudre dans $\R$ : (a) $x^4-13x^2+36=0$ ; (b) $x^4+3x^2-4=0$ ; (c) $x^4+x^2+1=0$.
\end{exo}
\begin{exo}[B]\label{c2e6}
Décomposer en éléments simples $\dfrac{2x-1}{(x-1)(x+3)}$ et $\dfrac{x^2+3x}{(x-1)(x+2)}$.
\end{exo}
\begin{exo}[C]\label{c2e7}
Décomposer $\dfrac1{(x-1)^2(x+1)}$ sous la forme $\frac A{x+1}+\frac B{x-1}+\frac C{(x-1)^2}$.
\end{exo}
\begin{exo}[C]\label{c2e8}
Montrer par $\varepsilon$--$\delta$ que $\lim_{x\to3}x^2=9$.
\end{exo}
\begin{exo}[C]\label{c2e9}
Montrer que $\lim_{x\to0}\frac{|x|}x$ n'existe pas (on raisonnera avec $\varepsilon=1$).
\end{exo}

\ifcorr
\subsection{Corrigés détaillés}
\begin{corrige}{\ref{c2e1}}
$80\times1{,}15\times0{,}85=78{,}2$ \texteuro{} ; facteur global $0{,}9775$, soit $-2{,}25\,\%$ (jamais $0\,\%$). Prix initial : $\frac{56}{0{,}8}=70$ \texteuro{}.
\end{corrige}
\begin{corrige}{\ref{c2e2}}
Pente $\frac{4-12}{40}=-0{,}2$ ; $T(t)=12-0{,}2t$ ; $T(25)=7^\circ$C ; $T=0$ pour $t=60$ min. Le résultat à $t=60$ est une \emph{extrapolation} hors de $[0;40]$ : elle vaut sous l'hypothèse d'un refroidissement affine, peu réaliste (une loi de refroidissement réelle est exponentielle).
\end{corrige}
\begin{corrige}{\ref{c2e3}}
$\Delta=25+24=49$, $x=\frac{5\pm7}4$ : $x=3$ ou $x=-\frac12$. Factorisation : $2(x-3)(x+\frac12)=(x-3)(2x+1)$. Contrôle : somme $\frac52=-\frac ba$, produit $-\frac32=\frac ca$. $\checkmark$
\end{corrige}
\begin{corrige}{\ref{c2e4}}
Le coefficient de $x^2$ est $1\ne0$ pour tout $m$ : pas de cas dégénéré. $\Delta=4m^2-4(m+2)=4(m^2-m-2)=4(m-2)(m+1)$. Donc : $m<-1$ ou $m>2$ : deux solutions ; $m=-1$ ou $m=2$ : une solution double ($x=m$) ; $-1<m<2$ : aucune. Contrôle $m=2$ : $x^2-4x+4=(x-2)^2$. $\checkmark$ ; $m=0$ : $x^2+2=0$ sans solution réelle. $\checkmark$
\end{corrige}
\begin{corrige}{\ref{c2e5}}
(a) $X^2-13X+36=0$, $X\in\{4,9\}$ : $x\in\{-3,-2,2,3\}$. (b) $X^2+3X-4=(X+4)(X-1)$, $X\in\{1,-4\}$ : seule $X=1$ est admissible, $x=\pm1$. (c) $X^2+X+1=0$, $\Delta=-3<0$ : aucune racine $X$ réelle, donc aucune solution réelle.
\end{corrige}
\begin{corrige}{\ref{c2e6}}
Première fraction : $A=\frac{2-1}{1+3}=\frac14$, $B=\frac{-6-1}{-3-1}=\frac74$ ; contrôle en $x=0$ : $\frac{-1}{-3}=\frac13$ et $-\frac14+\frac7{12}=\frac13$. Seconde : $\deg$ égaux, donc $\frac{x^2+3x}{x^2+x-2}=1+\frac{2x+2}{(x-1)(x+2)}$ ; $A=\frac{4}{3}$, $B=\frac{-2}{-3}=\frac23$ ; résultat $1+\frac4{3(x-1)}+\frac2{3(x+2)}$ ; contrôle en $x=0$ : $0$ et $1-\frac43+\frac13=0$. $\checkmark$
\end{corrige}
\begin{corrige}{\ref{c2e7}}
$C=\frac1{x+1}\big|_{x=1}=\frac12$ ; $A=\frac1{(x-1)^2}\big|_{x=-1}=\frac14$ ; en multipliant par $x$ et en faisant $x\to\infty$ : $0=A+B$, donc $B=-\frac14$. Résultat : $\frac1{4(x+1)}-\frac1{4(x-1)}+\frac1{2(x-1)^2}$. Contrôle en $x=0$ : $\frac1{(1)(1)}=1$ et $\frac14+\frac14+\frac12=1$. $\checkmark$
\end{corrige}
\begin{corrige}{\ref{c2e8}}
$|x^2-9|=|x-3||x+3|$. Si $|x-3|<1$ alors $x\in]2;4[$ et $|x+3|<7$. Avec $\delta=\min(1,\frac\varepsilon7)$ : $|x^2-9|<7\cdot\frac\varepsilon7=\varepsilon$.
\end{corrige}
\begin{corrige}{\ref{c2e9}}
Par l'absurde, supposons $\lim=L$. Pour $\varepsilon=1$ il existe $\delta>0$ tel que $0<|x|<\delta\Rightarrow\big|\frac{|x|}x-L\big|<1$. Avec $x=\frac\delta2$ : $|1-L|<1$ ; avec $x=-\frac\delta2$ : $|-1-L|<1$. L'inégalité triangulaire donne $2=|1-(-1)|\le|1-L|+|L+1|<2$, contradiction.
\end{corrige}
\else\stepcounter{subsection}\fi

\subsection{Contrôle de fin de semaine (30 minutes, 20 points)}
\begin{enumerate}
 \item \bareme{6} Un prix subit une hausse de $20\,\%$ puis une baisse de $10\,\%$. Donner la variation globale, puis le prix initial sachant que le prix final est $108$ \texteuro.
 \item \bareme{7} Résoudre dans $\R$ : $x^4-13x^2+36=0$.
 \item \bareme{7} Décomposer $\dfrac{5x+1}{(x-2)(x+1)}$ et vérifier le résultat en $x=0$.
\end{enumerate}
\ifcorr
\textbf{Corrigé et barème.} \textbf{1.} Facteur $1{,}2\times0{,}9=1{,}08$ : $+8\,\%$ \bareme{3} ; prix initial $\frac{108}{1{,}08}=100$ \texteuro{} \bareme{3}. \textbf{2.} $X=4$ ou $9$ \bareme{3} ; $x\in\{-3,-2,2,3\}$ \bareme{4}. \textbf{3.} $A=\frac{11}3$, $B=\frac43$ \bareme{4} ; contrôle en $0$ : $-\frac12$ et $-\frac{11}6+\frac43=-\frac12$ \bareme{3}. Seuil conseillé : $14/20$.
\fi

\subsection{Grille d'auto-évaluation}
$\square$ je chaîne des pourcentages par des facteurs ; $\square$ je teste le cas $a=0$ avant le discriminant ; $\square$ je rejette $X<0$ en bicarrée ; $\square$ je fais la division euclidienne quand $\deg P\ge$ degré du dénominateur ; $\square$ je choisis $\delta$ en fonction de $\varepsilon$ et jamais l'inverse.
\clearpage
\section{Dérivation, développements limités, optimisation}
\renewcommand{\titrecourant}{Semaine 3}

\begin{boussole}
\textbf{Objectif de la semaine.} Dériver sans hésiter, interpréter $f'$ (tangente, variations), approcher une fonction près d'un point par un polynôme, et résoudre un problème d'optimisation en justifiant que l'extremum trouvé est bien le bon.\\
\textbf{À la fin, je sais :}
\begin{enumerate}
 \item écrire l'équation de la tangente et reconnaître un point de non-dérivabilité ;
 \item établir le tableau de variations et classer les points critiques, y compris quand $f''=0$ ;
 \item écrire un développement de Taylor--Young d'ordre 2 et l'utiliser pour un calcul approché ou une limite ;
 \item modéliser un problème d'optimisation : variable, domaine, fonction, étude, conclusion ;
 \item démontrer le théorème de Fermat et la dérivée d'un produit.
\end{enumerate}
\textbf{Temps estimé} : 5 à 6 heures.
\end{boussole}

\subsection{Test d'entrée (5 minutes)}
\begin{enumerate}
 \item Dériver $x^3$, $\e^{2x}$, $\ln x$, $\sin x$.
 \item Quelle est la pente de la droite $y=3x-2$ ?
 \item Résoudre $x^2-4x+3\ge0$.
 \item Que vaut $\e^0$ ? $\ln 1$ ?
 \item Si $f'>0$ sur un intervalle, que dire de $f$ ?
\end{enumerate}
\textit{Réponses : $3x^2$, $2\e^{2x}$, $\frac1x$, $\cos x$ ; $3$ ; $]-\infty;1]\cup[3;+\infty[$ ; $1$ et $0$ ; $f$ est croissante.}

\subsection{Cours}
\begin{definition}[Dérivée]
$f$ est dérivable en $a$ si $\displaystyle\lim_{h\to0}\frac{f(a+h)-f(a)}h$ existe et est finie ; cette limite est $f'(a)$. La tangente en $a$ a pour équation $y=f(a)+f'(a)(x-a)$.
\end{definition}
\begin{propriete}[Règles]
Si $f,g$ sont dérivables : $(fg)'=f'g+fg'$ ; $\big(\frac fg\big)'=\frac{f'g-fg'}{g^2}$ (où $g\ne0$) ; $(f\circ g)'=(f'\circ g)\,g'$. Une fonction dérivable en $a$ est continue en $a$ ; la réciproque est fausse.
\end{propriete}
\begin{demonstration}[produit]
$f(x+h)g(x+h)-f(x)g(x)=f(x+h)\,[g(x+h)-g(x)]+g(x)\,[f(x+h)-f(x)]$. Divisons par $h$ et faisons $h\to0$ ; $f(x+h)\to f(x)$ car $f$ est continue (étant dérivable), d'où $f(x)g'(x)+g(x)f'(x)$.
\end{demonstration}
\begin{propriete}[Théorème de Fermat]
Si $f$ est dérivable en un point $a$ \emph{intérieur} à son intervalle de définition et admet en $a$ un extremum local, alors $f'(a)=0$.
\end{propriete}
\begin{demonstration}
Supposons un maximum local. Pour $h>0$ petit, $\frac{f(a+h)-f(a)}h\le0$, donc $f'(a)\le0$ par passage à la limite. Pour $h<0$ petit, le quotient est $\ge0$, donc $f'(a)\ge0$. Ainsi $f'(a)=0$. Le minimum se traite avec $-f$.
\end{demonstration}
\begin{piege}
La réciproque de Fermat est fausse (voir $x^3$) ; et le théorème ne dit rien aux \emph{bornes} du domaine ni aux points de non-dérivabilité.
\end{piege}
\begin{propriete}[Signe de $f'$ et variations]
Sur un intervalle, $f'\ge0$ implique $f$ croissante ; $f'\le0$ implique décroissante. Si $f'$ change de signe en $a$ (de $-$ à $+$), $f$ a un minimum local en $a$ ; de $+$ à $-$, un maximum local.
\end{propriete}
\begin{propriete}[Taylor--Young, ordre 2]
Si $f$ est deux fois dérivable en $a$, alors, quand $h\to0$ : $f(a+h)=f(a)+f'(a)h+\dfrac{f''(a)}2h^2+o(h^2)$.
\end{propriete}
Développements usuels en $0$ à l'ordre 2 : $\e^x=1+x+\frac{x^2}2+o(x^2)$ ; $\ln(1+x)=x-\frac{x^2}2+o(x^2)$ ; $\sqrt{1+x}=1+\frac x2-\frac{x^2}8+o(x^2)$ ; $\cos x=1-\frac{x^2}2+o(x^2)$ ; $\sin x=x+o(x^2)$.
\begin{methode}
\textbf{Optimisation.} (1) Choisir la variable et écrire son \emph{domaine} réaliste. (2) Exprimer la grandeur à optimiser en fonction d'elle seule. (3) Dériver, trouver les points critiques \emph{dans le domaine}. (4) Étudier le signe de $f'$ (ou comparer aux bornes). (5) Conclure par une phrase, avec unité, et contrôler par une valeur voisine.
\end{methode}

\subsection{Exemples travaillés et cas limites}
\begin{exemple}[tangente]
Tangente à $f(x)=x\e^x$ en $a=1$.\\
\emph{Solution.} $f'(x)=(1+x)\e^x$, donc $f(1)=\e$ et $f'(1)=2\e$ : $y=\e+2\e(x-1)=2\e x-\e$.
\end{exemple}
\begin{cas}[non-dérivabilité]
(a) $f(x)=|x|$ en $0$ : le quotient vaut $\pm1$ selon le signe de $h$, donc pas de limite ; $f$ est continue sans être dérivable. (b) $f(x)=x^{1/3}$ en $0$ : le quotient $h^{-2/3}\to+\infty$, la tangente est verticale, $f$ n'est pas dérivable. (c) $f(x)=x^2\sin\frac1x$, $f(0)=0$ : le quotient vaut $h\sin\frac1h\to0$, donc $f'(0)=0$ ; pourtant $f'(x)=2x\sin\frac1x-\cos\frac1x$ n'a pas de limite en $0$ : $f'$ n'est pas continue.
\end{cas}
\begin{exemple}[variations]
Étudier $f(x)=x^3-3x+1$.\\
\emph{Solution.} $f'(x)=3(x-1)(x+1)$. $f$ croît sur $]-\infty;-1]$, décroît sur $[-1;1]$, croît sur $[1;+\infty[$. Maximum local $f(-1)=3$, minimum local $f(1)=-1$. Comme $f(-1)>0>f(1)$, la courbe coupe l'axe des abscisses trois fois (trois racines réelles).
\end{exemple}
\begin{cas}[$f'(a)=0$ sans conclusion immédiate]
$x^3$ : $f'(0)=0$ mais $f'\ge0$ des deux côtés, pas d'extremum. $x^4$ : $f'(0)=f''(0)=0$, et pourtant minimum en $0$ (car $f'$ passe de $-$ à $+$). $|x|$ : minimum en $0$ sans que $f'(0)$ existe. Le critère fiable est le \emph{changement de signe de $f'$}, non la seule annulation.
\end{cas}
\begin{exemple}[Taylor--Young]
Valeur approchée de $\sqrt{1{,}02}$.\\
\emph{Solution.} Avec $x=0{,}02$ : $1+0{,}01-\frac{0{,}0004}8=1{,}00995$. Valeur exacte (calculée) : $1{,}00995049\ldots$ ; erreur $\approx5\times10^{-7}$.
\end{exemple}
\begin{cas}[limite de validité]
Un développement n'est valable que \emph{près} de $a$ : pour $\sqrt{1+x}$ avec $x=3$, l'ordre 2 donne $1+1{,}5-1{,}125=1{,}375$ au lieu de $2$. Et le terme $o(x^2)$ ne dit pas \og petit \fg{} en valeur absolue, seulement négligeable devant $x^2$ quand $x\to0$.
\end{cas}
\begin{exemple}[optimisation : boîte ouverte]
Dans une feuille carrée de $30$ cm de côté, on découpe aux quatre coins des carrés de côté $x$ pour former une boîte ouverte. Quel $x$ maximise le volume ?\\
\emph{Solution.} Domaine : $0<x<15$. $V(x)=x(30-2x)^2$, $V'(x)=(30-2x)(30-6x)=12(15-x)(5-x)$. $V'>0$ sur $]0;5[$, $V'<0$ sur $]5;15[$. Maximum en $x=5$ cm : $V=5\times400=2000$ cm$^3$.
\end{exemple}
\begin{cas}[contrainte sur le domaine]
Si l'on impose $x\le3$ (matériau disponible), le point critique $x=5$ est \emph{hors domaine}. $V$ étant croissante sur $]0;3]$, le maximum est atteint à la \emph{borne} $x=3$ : $V=3\times24^2=1728$ cm$^3$, et $V'(3)\ne0$. Ici le maximum n'est pas un point critique.
\end{cas}

\subsection{Exercices}
\begin{exo}[A]\label{c3e1}
Dériver : $f(x)=x^2\e^{-x}$ ; $g(x)=\ln(1+x^2)$ ; $h(x)=\dfrac{\sin x}x$ sur $\R^*$.
\end{exo}
\begin{exo}[A]\label{c3e2}
Donner l'équation de la tangente à $\ln$ en $a=\e$, puis vérifier que cette tangente passe par l'origine.
\end{exo}
\begin{exo}[A]\label{c3e3}
Dresser le tableau de variations de $f(x)=x^3-6x^2+9x$ et donner ses extremums locaux.
\end{exo}
\begin{exo}[B]\label{c3e4}
Donner le développement de $\ln(1+x)$ à l'ordre 2 en $0$ et approcher $\ln(1{,}05)$. Comparer à la valeur exacte arrondie.
\end{exo}
\begin{exo}[B]\label{c3e5}
Calculer $\lim_{x\to0}\dfrac{\e^x-1-x}{x^2}$ et $\lim_{x\to0}\dfrac{\cos x-1}{x^2}$ à l'aide de développements.
\end{exo}
\begin{exo}[B]\label{c3e6}
Étudier la dérivabilité en $0$ de $f(x)=x|x|$ et calculer $f'$. $f'$ est-elle dérivable en $0$ ?
\end{exo}
\begin{exo}[C]\label{c3e7}
Une boîte cylindrique fermée doit contenir $1000$ cm$^3$. Quel rayon minimise la surface de métal ? Donner la hauteur correspondante et conclure sur la forme.
\end{exo}
\begin{exo}[C]\label{c3e8}
Discuter, selon $m$, le nombre de racines réelles de $x^3-3x+m=0$.
\end{exo}
\begin{exo}[C]\label{c3e9}
Montrer qu'une fonction dérivable en $a$ est continue en $a$, et donner un exemple prouvant que la réciproque est fausse.
\end{exo}

\ifcorr
\subsection{Corrigés détaillés}
\begin{corrige}{\ref{c3e1}}
$f'(x)=2x\e^{-x}-x^2\e^{-x}=x(2-x)\e^{-x}$. $g'(x)=\frac{2x}{1+x^2}$ (dérivée de $\ln u$ : $\frac{u'}u$). $h'(x)=\frac{x\cos x-\sin x}{x^2}$. Contrôle de $f'$ en $x=2$ : $0$, cohérent avec le maximum local de $x^2\e^{-x}$ en $2$.
\end{corrige}
\begin{corrige}{\ref{c3e2}}
$f(\e)=1$, $f'(\e)=\frac1\e$ : $y=1+\frac{x-\e}\e=\frac x\e$. En $x=0$ : $y=0$, donc la tangente passe par l'origine.
\end{corrige}
\begin{corrige}{\ref{c3e3}}
$f'(x)=3x^2-12x+9=3(x-1)(x-3)$. $f$ croît sur $]-\infty;1]$, décroît sur $[1;3]$, croît sur $[3;+\infty[$. Maximum local $f(1)=4$, minimum local $f(3)=0$. Cohérence : $f(x)=x(x-3)^2$ s'annule en $3$ (racine double) et en $0$.
\end{corrige}
\begin{corrige}{\ref{c3e4}}
$\ln(1+x)=x-\frac{x^2}2+o(x^2)$. Avec $x=0{,}05$ : $0{,}05-0{,}00125=0{,}04875$. Valeur exacte arrondie : $0{,}04879$. Écart $\approx4\times10^{-5}$, du bon ordre de grandeur pour un terme d'ordre 3 ($\frac{x^3}3\approx4{,}2\times10^{-5}$).
\end{corrige}
\begin{corrige}{\ref{c3e5}}
$\e^x-1-x=\frac{x^2}2+o(x^2)$, donc la limite vaut $\frac12$. $\cos x-1=-\frac{x^2}2+o(x^2)$, donc la limite vaut $-\frac12$.
\end{corrige}
\begin{corrige}{\ref{c3e6}}
Le quotient en $0$ vaut $\frac{h|h|}h=|h|\to0$, donc $f'(0)=0$. Pour $x\ne0$ : $f'(x)=2|x|$ (pour $x>0$, $f=x^2$ ; pour $x<0$, $f=-x^2$). Ainsi $f'(x)=2|x|$ sur $\R$. Mais $f'=2|x|$ n'est pas dérivable en $0$ : $f$ est dérivable une fois, pas deux (pas de $f''(0)$).
\end{corrige}
\begin{corrige}{\ref{c3e7}}
Variable $r>0$, hauteur $h=\frac{1000}{\pi r^2}$. Surface $S(r)=2\pi r^2+2\pi rh=2\pi r^2+\frac{2000}r$. $S'(r)=4\pi r-\frac{2000}{r^2}=0\iff r^3=\frac{500}\pi$, $r\approx5{,}42$ cm. $S'<0$ avant, $>0$ après : minimum. $h\approx10{,}84\approx2r$, $S\approx553{,}6$ cm$^2$. Conclusion : la boîte optimale a une hauteur égale à son diamètre. Contrôle : avec $r=5$, $S=157{,}1+400=557{,}1>553{,}6$. $\checkmark$
\end{corrige}
\begin{corrige}{\ref{c3e8}}
$f(x)=x^3-3x+m$, $f'=3(x^2-1)$ : maximum local $f(-1)=m+2$, minimum local $f(1)=m-2$. Si $-2<m<2$ : $f(-1)>0>f(1)$, trois racines. Si $m=2$ : $f(-1)=4$, $f(1)=0$ : racine double en $1$ et racine simple (factorisation $(x-1)^2(x+2)$), donc deux racines distinctes ; de même $m=-2$ (symétrie $x\to-x$). Si $|m|>2$ : une seule racine.
\end{corrige}
\begin{corrige}{\ref{c3e9}}
Si $f$ est dérivable en $a$, $f(a+h)-f(a)=h\cdot\frac{f(a+h)-f(a)}h\to0\cdot f'(a)=0$, donc $f$ est continue en $a$. Réciproque fausse : $|x|$ en $0$ (continue, non dérivable).
\end{corrige}
\else\stepcounter{subsection}\fi

\subsection{Contrôle de fin de semaine (30 minutes, 20 points)}
\begin{enumerate}
 \item \bareme{6} Dériver $f(x)=x\e^{2x}$ et donner la tangente en $0$.
 \item \bareme{7} Étudier les variations de $f(x)=x^3-12x+1$ ; donner les extremums locaux.
 \item \bareme{7} Écrire le développement de $\sqrt{1+x}$ à l'ordre 2 et approcher $\sqrt{1{,}04}$.
\end{enumerate}
\ifcorr
\textbf{Corrigé et barème.} \textbf{1.} $f'(x)=(1+2x)\e^{2x}$ \bareme{3} ; $f(0)=0$, $f'(0)=1$, tangente $y=x$ \bareme{3}. \textbf{2.} $f'=3(x-2)(x+2)$ \bareme{3} ; croissante, décroissante, croissante \bareme{2} ; maximum $f(-2)=17$, minimum $f(2)=-15$ \bareme{2}. \textbf{3.} $\sqrt{1+x}=1+\frac x2-\frac{x^2}8+o(x^2)$ \bareme{4} ; $1+0{,}02-0{,}0002=1{,}0198$ (exact $1{,}0198039$) \bareme{3}. Seuil conseillé : $14/20$.
\fi

\subsection{Grille d'auto-évaluation}
$\square$ je donne la tangente sans me tromper de point ; $\square$ je cherche le changement de signe de $f'$ et pas seulement $f'=0$ ; $\square$ je vérifie que le point critique est dans le domaine ; $\square$ je compare aux bornes quand le domaine est fermé ; $\square$ je n'utilise Taylor--Young que près du point.
\clearpage
\section{Intégration}
\renewcommand{\titrecourant}{Semaine 4}

\begin{boussole}
\textbf{Objectif de la semaine.} Calculer une intégrale par primitive, par parties ou par changement de variable, en gérant les bornes ; décider si une intégrale impropre converge ; calculer une aire (avec croisements de courbes) et un centre de gravité.\\
\textbf{À la fin, je sais :}
\begin{enumerate}
 \item interpréter $\int_a^bf$ comme limite de sommes de Riemann ;
 \item appliquer le théorème fondamental et l'intégration par parties ;
 \item changer de variable \emph{avec les bornes} ;
 \item reconnaître une intégrale impropre convergente ou divergente ;
 \item calculer une aire entre deux courbes sans me tromper de signe ;
 \item calculer une valeur moyenne et un centre de gravité plan.
\end{enumerate}
\textbf{Temps estimé} : 5 à 6 heures.
\end{boussole}

\subsection{Test d'entrée (5 minutes)}
\begin{enumerate}
 \item Une primitive de $x^2$ ? de $\frac1x$ ? de $\e^{2x}$ ?
 \item Dériver $x\ln x-x$.
 \item Calculer $\int_0^1(2x+1)\,\dd x$.
 \item Quel est le signe de $\int_a^bf$ si $f\le0$ et $a<b$ ?
 \item Que vaut $\int_a^af$ ?
\end{enumerate}
\textit{Réponses : $\frac{x^3}3$, $\ln|x|$, $\frac{\e^{2x}}2$ ; $\ln x$ ; $2$ ; négatif ou nul ; $0$.}

\subsection{Cours}
\begin{definition}[Intégrale de Riemann]
Pour $f$ continue sur $[a;b]$, en subdivisant en $n$ pas égaux $\Delta x=\frac{b-a}n$, la somme $S_n=\sum_{k=0}^{n-1}f(a+k\Delta x)\Delta x$ converge quand $n\to\infty$ vers $\int_a^bf(x)\,\dd x$ (aire algébrique entre la courbe et l'axe).
\end{definition}
\begin{propriete}[Théorème fondamental]
Si $f$ est continue sur $[a;b]$ et $F$ est une primitive de $f$, alors $\int_a^bf=F(b)-F(a)$. De plus $x\mapsto\int_a^xf$ est dérivable de dérivée $f$.
\end{propriete}
\begin{propriete}[Intégration par parties]
Si $u,v$ sont de classe $C^1$ : $\displaystyle\int_a^bu\,v'=\big[uv\big]_a^b-\int_a^bu'\,v$.
\end{propriete}
\begin{demonstration}
Intégrer $(uv)'=u'v+uv'$ sur $[a;b]$ et appliquer le théorème fondamental.
\end{demonstration}
\begin{propriete}[Changement de variable]
Si $\varphi$ est $C^1$ et $f$ continue : $\displaystyle\int_a^bf(\varphi(t))\varphi'(t)\,\dd t=\int_{\varphi(a)}^{\varphi(b)}f(u)\,\dd u$. Les bornes \emph{changent avec la variable}.
\end{propriete}
\begin{propriete}[Valeur moyenne]
Si $f$ est continue sur $[a;b]$, $\mu=\frac1{b-a}\int_a^bf$ vérifie $\mu=f(c)$ pour au moins un $c\in[a;b]$.
\end{propriete}
\begin{definition}[Intégrale impropre]
$\int_a^{+\infty}f=\lim_{X\to+\infty}\int_a^Xf$ si cette limite est finie (convergence). Pour $f$ non bornée en $a$ : $\int_a^bf=\lim_{\eta\to0^+}\int_{a+\eta}^bf$.
\end{definition}
\begin{propriete}[Référence de Riemann]
$\int_1^{+\infty}x^{-p}\,\dd x$ converge si et seulement si $p>1$ ; $\int_0^1x^{-p}\,\dd x$ converge si et seulement si $p<1$.
\end{propriete}
\begin{definition}[Centre de gravité plan]
Pour la région $0\le y\le f(x)$, $a\le x\le b$, d'aire $A=\int_a^bf$ : $x_G=\frac1A\int_a^bxf(x)\,\dd x$ et $y_G=\frac1{2A}\int_a^bf(x)^2\,\dd x$.
\end{definition}
\begin{methode}
\textbf{Choisir la technique.} Primitive visible $\to$ théorème fondamental. Produit polynôme $\times$ $\e^x$, $\sin$, $\cos$ $\to$ parties (le polynôme joue $u$). $\ln x$ seul $\to$ parties avec $v'=1$. Un facteur est la dérivée d'une expression qui apparaît ailleurs ($\frac{2x}{1+x^2}$) $\to$ changement de variable.
\end{methode}
\begin{piege}
(i) Oublier de changer les bornes ou de revenir à $x$. (ii) Calculer une \og aire \fg{} avec une intégrale signée sans découper aux croisements. (iii) Appliquer le théorème fondamental à une fonction non continue sur $[a;b]$ (pôle à l'intérieur).
\end{piege}

\subsection{Exemples travaillés et cas limites}
\begin{exemple}[somme de Riemann]
Calculer $\int_0^1x^2\,\dd x$ par la définition.\\
\emph{Solution.} $S_n=\sum_{k=0}^{n-1}\frac{k^2}{n^2}\cdot\frac1n=\frac{(n-1)n(2n-1)}{6n^3}=\frac{(n-1)(2n-1)}{6n^2}\to\frac13$. Contrôle $n=4$ : $\frac{0+1+4+9}{64}=\frac7{32}$ ; formule : $\frac{3\cdot7}{96}=\frac7{32}$. $\checkmark$
\end{exemple}
\begin{exemple}[intégration par parties]
Calculer $\int_0^1x\e^x\,\dd x$.\\
\emph{Solution.} $u=x$, $v'=\e^x$ : $[x\e^x]_0^1-\int_0^1\e^x=\e-(\e-1)=1$. Contrôle : $(x-1)\e^x$ est une primitive, valeurs $0$ en $1$ et $-1$ en $0$. $\checkmark$
\end{exemple}
\begin{cas}[parties cycliques et $\ln$]
(a) $\int\ln x\,\dd x$ : $u=\ln x$, $v'=1$ donne $x\ln x-x$ (une primitive de $\ln$, par parties alors qu'il n'y a pas de produit apparent). (b) $J=\int_0^\pi\e^x\sin x\,\dd x$ : deux parties successives redonnent $J$ au second membre : $J=\e^\pi+1-J$, donc $2J=\e^\pi+1$ et $J=\frac{\e^\pi+1}2$. On résout une équation en $J$ au lieu de continuer indéfiniment.
\end{cas}
\begin{exemple}[changement de variable]
Calculer $\int_0^1\frac{2x}{1+x^2}\,\dd x$.\\
\emph{Solution.} $u=1+x^2$, $\dd u=2x\,\dd x$ ; bornes $u:1\to2$ : $\int_1^2\frac{\dd u}u=\ln2$.
\end{exemple}
\begin{cas}[symétrie et pièges de bornes]
$\int_{-1}^1x\e^{x^2}\,\dd x=0$ : l'intégrande est impaire sur un intervalle symétrique (on le confirme avec $u=x^2$, bornes $1\to1$, intégrale nulle). Ne pas écrire les bornes $-1$ et $1$ dans $u$ puis oublier que $u(-1)=u(1)$.
\end{cas}
\begin{exemple}[impropre convergente]
$\int_0^1x^{-1/2}\,\dd x=\lim_{\eta\to0^+}[2\sqrt x]_\eta^1=2$ : l'aire sous une courbe non bornée peut être finie.
\end{exemple}
\begin{cas}[impropres divergentes]
$\int_0^1\frac{\dd x}x=\lim[\ln x]_\eta^1=+\infty$ : divergente. Plus insidieux : $\int_{-1}^1\frac{\dd x}{x^2}$ ne vaut \emph{pas} $[-\frac1x]_{-1}^1=-2$ : l'intégrande est positive, l'intégrale vaut $+\infty$ car $\frac1{x^2}$ n'est pas continue en $0\in[-1;1]$. Le théorème fondamental ne s'applique pas à travers le pôle. Enfin $\int_1^{+\infty}\frac{\dd x}x$ diverge alors que $\int_1^{+\infty}\frac{\dd x}{x^2}=1$.
\end{cas}
\begin{exemple}[aire entre deux courbes]
Aire comprise entre $y=x$ et $y=x^2$ sur $[0;1]$ : $\int_0^1(x-x^2)\,\dd x=\frac12-\frac13=\frac16$ (la courbe du dessus est $y=x$).
\end{exemple}
\begin{cas}[croisement : l'intégrale n'est pas l'aire]
Aire entre $y=x^2-2x$ et l'axe des abscisses sur $[0;3]$. L'expression change de signe en $x=2$. L'intégrale signée $\int_0^3(x^2-2x)\,\dd x=0$, alors que l'aire vaut $\int_0^2(2x-x^2)+\int_2^3(x^2-2x)=\frac43+\frac43=\frac83$. Toujours chercher les changements de signe avant d'intégrer.
\end{cas}
\begin{exemple}[centre de gravité]
Région sous $y=4-x^2$ sur $[0;2]$.\\
\emph{Solution.} $A=\int_0^2(4-x^2)=\frac{16}3$ ; $x_G=\frac1A\int_0^2(4x-x^3)=\frac{4}{16/3}=\frac34$ ; $y_G=\frac1{2A}\int_0^2(4-x^2)^2=\frac{3}{32}\cdot\frac{256}{15}=\frac85$. Contrôle : $y_G<4$ et $x_G<1$ (région plus massive près de l'axe $Oy$). $\checkmark$
\end{exemple}

\subsection{Exercices}
\begin{exo}[A]\label{c4e1}
Calculer $\int_0^2(3x^2-2x+1)\,\dd x$ et $\int_0^{\pi/2}\sin x\cos x\,\dd x$.
\end{exo}
\begin{exo}[A]\label{c4e2}
Calculer $\int_0^1x\sqrt{1+x^2}\,\dd x$ par changement de variable.
\end{exo}
\begin{exo}[A]\label{c4e3}
Calculer $\int_1^{\e}\ln x\,\dd x$.
\end{exo}
\begin{exo}[B]\label{c4e4}
Calculer $\int_0^\pi\e^x\sin x\,\dd x$.
\end{exo}
\begin{exo}[B]\label{c4e5}
Calculer l'aire comprise entre la courbe de $f(x)=x^3-x$ et l'axe des abscisses sur $[-1;1]$. Comparer à $\int_{-1}^1f$.
\end{exo}
\begin{exo}[B]\label{c4e6}
Déterminer la valeur moyenne de $f(x)=x^2$ sur $[0;1]$ et le point $c$ prévu par la propriété de valeur moyenne.
\end{exo}
\begin{exo}[C]\label{c4e7}
Pour quels $p>0$ l'intégrale $\int_1^{+\infty}x^{-p}\,\dd x$ converge-t-elle ? Même question pour $\int_0^1x^{-p}\,\dd x$.
\end{exo}
\begin{exo}[C]\label{c4e8}
Déterminer le centre de gravité de la région sous $y=x^2$ sur $[0;1]$.
\end{exo}
\begin{exo}[C]\label{c4e9}
Montrer que $\int_0^{+\infty}x\e^{-x}\,\dd x$ converge et calculer sa valeur.
\end{exo}

\ifcorr
\subsection{Corrigés détaillés}
\begin{corrige}{\ref{c4e1}}
$[x^3-x^2+x]_0^2=8-4+2=6$. Pour la seconde, $\sin x\cos x=\frac12\sin2x$ : $\big[-\frac14\cos2x\big]_0^{\pi/2}=\frac14+\frac14=\frac12$ (ou $u=\sin x$ : $\int_0^1u\,\dd u=\frac12$).
\end{corrige}
\begin{corrige}{\ref{c4e2}}
$u=1+x^2$, $\dd u=2x\,\dd x$, bornes $1\to2$ : $\frac12\int_1^2\sqrt u\,\dd u=\frac12\cdot\frac23[u^{3/2}]_1^2=\frac13(2\sqrt2-1)\approx0{,}609$. Contrôle : intégrande comprise entre $0$ et $\sqrt2$, intégrale plausible. $\checkmark$
\end{corrige}
\begin{corrige}{\ref{c4e3}}
$[x\ln x-x]_1^{\e}=(\e-\e)-(0-1)=1$.
\end{corrige}
\begin{corrige}{\ref{c4e4}}
$J=[\e^x\sin x]_0^\pi-\int_0^\pi\e^x\cos x=-\int_0^\pi\e^x\cos x$ ; seconde parties : $\int_0^\pi\e^x\cos x=[\e^x\cos x]_0^\pi+J=-\e^\pi-1+J$. Donc $J=\e^\pi+1-J$, soit $J=\frac{\e^\pi+1}2\approx12{,}07$. Contrôle par primitive $\frac{\e^x(\sin x-\cos x)}2$. $\checkmark$
\end{corrige}
\begin{corrige}{\ref{c4e5}}
$f(x)=x(x-1)(x+1)$ : positive sur $]-1;0[$, négative sur $]0;1[$. $\int_{-1}^1f=0$ (impaire). Aire : $2\int_0^1(x-x^3)\,\dd x=2\big(\frac12-\frac14\big)=\frac12$.
\end{corrige}
\begin{corrige}{\ref{c4e6}}
$\mu=\int_0^1x^2=\frac13$. $c^2=\frac13\Rightarrow c=\frac1{\sqrt3}\approx0{,}577\in[0;1]$.
\end{corrige}
\begin{corrige}{\ref{c4e7}}
$p\ne1$ : $\int_1^X x^{-p}=\frac{X^{1-p}-1}{1-p}$, limite finie si et seulement si $1-p<0$, c'est-à-dire $p>1$ ; $p=1$ : $\ln X\to\infty$. En $0$ : $\int_\eta^1x^{-p}=\frac{1-\eta^{1-p}}{1-p}$, finie si et seulement si $1-p>0$, soit $p<1$ ; $p=1$ : diverge. Conclusion : la valeur $p=1$ diverge dans les deux cas.
\end{corrige}
\begin{corrige}{\ref{c4e8}}
$A=\frac13$ ; $x_G=3\int_0^1x^3=\frac34$ ; $y_G=\frac{1}{2A}\int_0^1x^4=\frac32\cdot\frac15=\frac3{10}$. Contrôle : $x_G>\frac12$ car la région est plus large vers $x=1$, et $y_G<\frac12$ car la courbe est basse sur la majeure partie de l'intervalle. $\checkmark$
\end{corrige}
\begin{corrige}{\ref{c4e9}}
Parties : $\int_0^Xx\e^{-x}=[-x\e^{-x}]_0^X+\int_0^X\e^{-x}=-X\e^{-X}+1-\e^{-X}\to1$ (croissances comparées). Convergente de valeur $1$.
\end{corrige}
\else\stepcounter{subsection}\fi

\subsection{Contrôle de fin de semaine (30 minutes, 20 points)}
\begin{enumerate}
 \item \bareme{7} Calculer $\int_1^{\e}x\ln x\,\dd x$.
 \item \bareme{6} Calculer $\int_0^1\frac x{1+x^2}\,\dd x$ par changement de variable.
 \item \bareme{7} Étudier la convergence de $\int_0^1\frac{\dd x}{\sqrt x}$ et de $\int_0^1\frac{\dd x}x$, avec la valeur quand elle existe.
\end{enumerate}
\ifcorr
\textbf{Corrigé et barème.} \textbf{1.} $u=\ln x$, $v'=x$ : $\big[\frac{x^2}2\ln x\big]_1^{\e}-\int_1^{\e}\frac x2=\frac{\e^2}2-\frac{\e^2-1}4=\frac{\e^2+1}4$ \bareme{3 pour la méthode, 4 pour le résultat}. \textbf{2.} $u=1+x^2$, $\frac12\int_1^2\frac{\dd u}u=\frac{\ln2}2$ \bareme{3+3}. \textbf{3.} Première : $2$, convergente \bareme{4} ; seconde : divergente ($\ln$ tend vers $-\infty$ en $0$) \bareme{3}. Seuil conseillé : $14/20$.
\fi

\subsection{Grille d'auto-évaluation}
$\square$ je change les bornes avec la variable ; $\square$ je découpe l'aire aux changements de signe ; $\square$ je vérifie la continuité avant d'appliquer le théorème fondamental ; $\square$ je mets une limite sur toute borne impropre ; $\square$ je contrôle par un ordre de grandeur.
\clearpage
\section{Coordonnées, produit scalaire, plans et volumes}
\renewcommand{\titrecourant}{Semaine 5}

\begin{boussole}
\textbf{Objectif de la semaine.} Passer d'un système de coordonnées à un autre sans se tromper d'angle, utiliser le produit scalaire (angle, projection, distance à un plan), écrire l'équation d'un plan, et calculer volumes et masses de solides usuels.\\
\textbf{Convention de ce cahier} \verif : en coordonnées sphériques, $\rho$ est la distance à l'origine, $\varphi\in[0;\pi]$ l'angle avec l'axe $Oz$ et $\theta$ l'angle d'azimut ; d'autres cours échangent les noms $\theta$ et $\varphi$.\\
\textbf{À la fin, je sais :}
\begin{enumerate}
 \item convertir polaires, cylindriques et sphériques en cartésiennes, et inversement, avec le bon quadrant ;
 \item calculer un angle, une projection orthogonale, une distance point--plan ;
 \item écrire l'équation d'un plan et reconnaître les cas dégénérés ;
 \item donner le volume d'un cylindre, d'un cône, d'une boule ; calculer une masse à densité variable.
\end{enumerate}
\textbf{Temps estimé} : 5 à 6 heures.
\end{boussole}

\subsection{Test d'entrée (5 minutes)}
\begin{enumerate}
 \item Théorème de Pythagore dans $\R^3$ : norme de $(1,2,2)$ ?
 \item Valeurs de $\cos\frac\pi3$, $\sin\frac\pi6$, $\cos\frac{2\pi}3$.
 \item Dans quel quadrant est le point $(-1,-1)$ ?
 \item Aire du disque de rayon $R$ ? Volume d'un cylindre ?
 \item Que vaut $\int_0^R r^2\,\dd r$ ?
\end{enumerate}
\textit{Réponses : $3$ ; $\frac12$, $\frac12$, $-\frac12$ ; troisième quadrant ; $\pi R^2$, $\pi R^2h$ ; $\frac{R^3}3$.}

\subsection{Cours}
\begin{definition}[Coordonnées]
\emph{Polaires} $(r,\theta)$ : $x=r\cos\theta$, $y=r\sin\theta$, $r\ge0$. \emph{Cylindriques} $(r,\theta,z)$ : polaires pour $(x,y)$ plus $z$. \emph{Sphériques} $(\rho,\varphi,\theta)$ : $x=\rho\sin\varphi\cos\theta$, $y=\rho\sin\varphi\sin\theta$, $z=\rho\cos\varphi$.
\end{definition}
\begin{propriete}[Passage inverse]
$r=\sqrt{x^2+y^2}$ ; $\rho=\sqrt{x^2+y^2+z^2}$ ; $\cos\varphi=\frac z\rho$ ($\rho\ne0$) ; $\theta$ est l'angle de $(x,y)$ avec l'axe $Ox$, déterminé par \emph{les signes de $x$ et de $y$} (jamais par $\arctan\frac yx$ seul).
\end{propriete}
\begin{definition}[Produit scalaire]
$u\cdot v=u_1v_1+u_2v_2+u_3v_3=\|u\|\|v\|\cos\gamma$ où $\gamma$ est l'angle géométrique entre $u$ et $v$ non nuls. Deux vecteurs sont orthogonaux si et seulement si $u\cdot v=0$.
\end{definition}
\begin{propriete}[Cauchy--Schwarz]
$|u\cdot v|\le\|u\|\|v\|$, avec égalité si et seulement si $u$ et $v$ sont colinéaires.
\end{propriete}
\begin{demonstration}
Si $v=0$ c'est clair. Sinon, pour tout $\lambda$, $0\le\|u-\lambda v\|^2=\|u\|^2-2\lambda\,u\cdot v+\lambda^2\|v\|^2$. Ce trinôme en $\lambda$ reste positif ou nul, donc son discriminant est $\le0$ : $4(u\cdot v)^2-4\|u\|^2\|v\|^2\le0$. L'égalité correspond à un discriminant nul, c'est-à-dire à un $\lambda$ avec $u=\lambda v$.
\end{demonstration}
\begin{propriete}[Projection et distance]
La projection orthogonale de $v$ sur la droite dirigée par $u\ne0$ est $\dfrac{u\cdot v}{u\cdot u}\,u$. La distance du point $M(x_0,y_0,z_0)$ au plan $ax+by+cz=d$ est $\dfrac{|ax_0+by_0+cz_0-d|}{\sqrt{a^2+b^2+c^2}}$.
\end{propriete}
\begin{propriete}[Équation d'un plan]
Le plan passant par $A$ de vecteur normal $n\ne0$ est $n\cdot(X-A)=0$. Si $n=(a,b,c)$, c'est $ax+by+cz=d$ avec $d=n\cdot A$.
\end{propriete}
\begin{propriete}[Volumes usuels]
Cylindre : $\pi R^2h$. Cône : $\frac13\pi R^2H$. Boule : $\frac43\pi R^3$. Élément de volume : $r\,\dd r\,\dd\theta\,\dd z$ (cylindriques), $\rho^2\sin\varphi\,\dd\rho\,\dd\varphi\,\dd\theta$ (sphériques).
\end{propriete}
\begin{methode}
\textbf{Convertir un angle.} (1) Calculer $r$. (2) Repérer le quadrant par les signes de $x$ et $y$. (3) Donner $\theta$ dans $[0;2\pi[$ (ou $]-\pi;\pi]$, en l'indiquant). (4) Contrôler en reconvertissant.
\end{methode}
\begin{piege}
(i) $\theta=\arctan\frac yx$ donne un angle dans $]-\frac\pi2;\frac\pi2[$ : faux dans deux quadrants sur quatre. (ii) Confondre les rôles de $\theta$ et $\varphi$ selon le cours. (iii) Oublier le facteur $r$ ou $\rho^2\sin\varphi$ dans l'élément de volume. (iv) Utiliser la formule de distance sans écrire le plan sous forme $ax+by+cz=d$.
\end{piege}

\subsection{Exemples travaillés et cas limites}
\begin{exemple}[quadrant]
Coordonnées polaires de $(-1,-1)$.\\
\emph{Solution.} $r=\sqrt2$ ; $x<0$, $y<0$ : troisième quadrant, $\theta=\frac{5\pi}4$ (et non $\arctan1=\frac\pi4$, qui désigne $(1,1)$).
\end{exemple}
\begin{cas}[points sur les axes, origine]
$(0,-2)$ : $r=2$, $\theta=\frac{3\pi}2$ (ou $-\frac\pi2$) ; le quotient $\frac yx$ est indéfini, la méthode \og arctan \fg{} échoue. L'origine : $r=0$ et $\theta$ est \emph{indéterminé} (tout angle convient) ; en sphériques, $\rho=0$ rend $\varphi$ et $\theta$ indéterminés. Sur l'axe $Oz$ ($x=y=0$, $z\ne0$) : $\theta$ est indéterminé, et $\varphi=0$ si $z>0$, $\varphi=\pi$ si $z<0$.
\end{cas}
\begin{exemple}[cartésiennes $\to$ sphériques]
$M(1,1,\sqrt2)$.\\
\emph{Solution.} $\rho=\sqrt{1+1+2}=2$ ; $r=\sqrt2$ ; $\cos\varphi=\frac{\sqrt2}2$, $\varphi=\frac\pi4$ ; $x,y>0$ : $\theta=\frac\pi4$. Contrôle : $x=2\sin\frac\pi4\cos\frac\pi4=1$, $z=2\cos\frac\pi4=\sqrt2$. $\checkmark$
\end{exemple}
\begin{cas}[courbe en polaires]
$r=2\cos\theta$ : en multipliant par $r$, $x^2+y^2=2x$, soit $(x-1)^2+y^2=1$ : le cercle de centre $(1,0)$, rayon $1$. Pour $\theta\in]\frac\pi2;\frac{3\pi}2[$, $r$ est négatif : avec la convention $r\ge0$ ces angles ne donnent rien, et le cercle est parcouru une seule fois pour $\theta\in[-\frac\pi2;\frac\pi2]$.
\end{cas}
\begin{exemple}[produit scalaire, projection, distance]
$u=(1,2,2)$, $v=(3,0,4)$.\\
\emph{Solution.} $u\cdot v=11$, $\|u\|=3$, $\|v\|=5$, $\cos\gamma=\frac{11}{15}$. Projection de $v$ sur $u$ : $\frac{11}9u=\big(\frac{11}9,\frac{22}9,\frac{22}9\big)$. Distance de $v$ au plan $x+2y+2z=0$ : $\frac{|3+0+8|}3=\frac{11}3$. Contrôle : $\|v\|\cos\gamma=\frac{11}3$, c'est la longueur algébrique de la projection. $\checkmark$
\end{exemple}
\begin{cas}[vecteurs nuls ou orthogonaux]
Si $u=0$, l'angle n'est pas défini et la formule de projection divise par $u\cdot u=0$. Si $u\cdot v=0$ avec $u,v\ne0$, les vecteurs sont orthogonaux : $\gamma=\frac\pi2$ et la projection est le vecteur nul. Si $u=\lambda v$ avec $\lambda<0$, on a $\gamma=\pi$ et Cauchy--Schwarz est une égalité, avec $u\cdot v=-\|u\|\|v\|<0$.
\end{cas}
\begin{exemple}[distance à un plan]
Plan $2x-y+2z=6$. Distance de l'origine : $\frac{|-6|}3=2$. Distance de $(1,1,1)$ : $\frac{|2-1+2-6|}3=1$.
\end{exemple}
\begin{exemple}[masse à densité variable]
Une boule de rayon $R$ a une densité $\rho_0\big(1-\frac rR\big)$ à la distance $r$ du centre. Masse ?\\
\emph{Solution.} En coquilles de volume $4\pi r^2\,\dd r$ : $M=4\pi\rho_0\int_0^R\big(r^2-\frac{r^3}R\big)\dd r=4\pi\rho_0\big(\frac{R^3}3-\frac{R^3}4\big)=\frac{\pi\rho_0R^3}3$. Contrôle : c'est le quart de la masse $\frac43\pi\rho_0R^3$ d'une boule homogène de densité $\rho_0$, ce qui est cohérent avec une densité qui décroît de $\rho_0$ à $0$.
\end{exemple}

\subsection{Exercices}
\begin{exo}[A]\label{c5e1}
Donner les coordonnées polaires de $(-\sqrt3,1)$ et les cartésiennes du point $(r=3,\theta=\frac{2\pi}3)$.
\end{exo}
\begin{exo}[A]\label{c5e2}
Un point a pour coordonnées sphériques $\rho=2$, $\varphi=\frac\pi6$, $\theta=\frac\pi3$. Donner ses coordonnées cartésiennes et contrôler $\rho$.
\end{exo}
\begin{exo}[A]\label{c5e3}
Pour $u=(2,-1,2)$ et $v=(1,2,2)$, calculer $u\cdot v$, les normes et $\cos\gamma$.
\end{exo}
\begin{exo}[B]\label{c5e4}
Quelle est la courbe d'équation polaire $r=4\sin\theta$ ? Donner son équation cartésienne.
\end{exo}
\begin{exo}[B]\label{c5e5}
Déterminer l'équation du plan passant par $A(1,0,0)$, $B(0,2,0)$, $C(0,0,3)$, puis sa distance à l'origine.
\end{exo}
\begin{exo}[B]\label{c5e6}
Calculer les volumes : cylindre $R=2$, $h=3$ ; cône $R=3$, $H=4$ ; boule $R=2$.
\end{exo}
\begin{exo}[C]\label{c5e7}
Volume du solide $\{\rho\le1,\ \varphi\le\frac\pi4\}$ (cône sphérique \og cornet de glace \fg{}).
\end{exo}
\begin{exo}[C]\label{c5e8}
Montrer l'inégalité triangulaire $\|u+v\|\le\|u\|+\|v\|$ à partir de Cauchy--Schwarz.
\end{exo}
\begin{exo}[C]\label{c5e9}
Masse de la boule de rayon $1$ de densité $\rho(r)=r$.
\end{exo}

\ifcorr
\subsection{Corrigés détaillés}
\begin{corrige}{\ref{c5e1}}
$(-\sqrt3,1)$ : $r=2$, $x<0$, $y>0$ (deuxième quadrant), $\cos\theta=-\frac{\sqrt3}2$, $\sin\theta=\frac12$ : $\theta=\frac{5\pi}6$. Pour le second : $x=3\cos\frac{2\pi}3=-\frac32$, $y=3\sin\frac{2\pi}3=\frac{3\sqrt3}2$.
\end{corrige}
\begin{corrige}{\ref{c5e2}}
$x=2\cdot\frac12\cdot\frac12=\frac12$ ; $y=2\cdot\frac12\cdot\frac{\sqrt3}2=\frac{\sqrt3}2$ ; $z=2\cos\frac\pi6=\sqrt3$. Contrôle : $\frac14+\frac34+3=4=\rho^2$. $\checkmark$
\end{corrige}
\begin{corrige}{\ref{c5e3}}
$u\cdot v=2-2+4=4$ ; $\|u\|=\|v\|=3$ ; $\cos\gamma=\frac49$, $\gamma\approx63{,}6^\circ$.
\end{corrige}
\begin{corrige}{\ref{c5e4}}
$r^2=4r\sin\theta$ donne $x^2+y^2=4y$, soit $x^2+(y-2)^2=4$ : cercle de centre $(0,2)$ et de rayon $2$, parcouru pour $\theta\in[0;\pi]$ (là où $r\ge0$).
\end{corrige}
\begin{corrige}{\ref{c5e5}}
$\overrightarrow{AB}=(-1,2,0)$, $\overrightarrow{AC}=(-1,0,3)$, normal $\overrightarrow{AB}\times\overrightarrow{AC}=(6,3,2)$. Plan : $6x+3y+2z=6$ (vérifié en $A$, $B$, $C$ : $6=6$). Distance à l'origine : $\frac{6}{\sqrt{36+9+4}}=\frac67$.
\end{corrige}
\begin{corrige}{\ref{c5e6}}
Cylindre : $\pi\cdot4\cdot3=12\pi$. Cône : $\frac13\pi\cdot9\cdot4=12\pi$. Boule : $\frac43\pi\cdot8=\frac{32\pi}3$.
\end{corrige}
\begin{corrige}{\ref{c5e7}}
$V=\int_0^{2\pi}\!\!\int_0^{\pi/4}\!\!\int_0^1\rho^2\sin\varphi\,\dd\rho\,\dd\varphi\,\dd\theta=2\pi\cdot\frac13\cdot\big(1-\cos\frac\pi4\big)=\frac{\pi(2-\sqrt2)}3\approx0{,}613$. Contrôle : fraction $\frac{1-\cos(\pi/4)}2\approx0{,}146$ de la boule unité $\frac{4\pi}3\approx4{,}19$, soit $0{,}613$. $\checkmark$
\end{corrige}
\begin{corrige}{\ref{c5e8}}
$\|u+v\|^2=\|u\|^2+2u\cdot v+\|v\|^2\le\|u\|^2+2\|u\|\|v\|+\|v\|^2=(\|u\|+\|v\|)^2$ car $u\cdot v\le|u\cdot v|\le\|u\|\|v\|$. On conclut par croissance de la racine carrée.
\end{corrige}
\begin{corrige}{\ref{c5e9}}
$M=\int_0^1r\cdot4\pi r^2\,\dd r=4\pi\cdot\frac14=\pi$.
\end{corrige}
\else\stepcounter{subsection}\fi

\subsection{Contrôle de fin de semaine (30 minutes, 20 points)}
\begin{enumerate}
 \item \bareme{6} Donner les coordonnées polaires de $(\sqrt3,-1)$ ($\theta\in]-\pi;\pi]$).
 \item \bareme{7} Pour $u=(1,-1,2)$ et $v=(2,0,1)$, calculer $u\cdot v$, $\cos\gamma$ et dire si $\gamma$ est aigu.
 \item \bareme{7} Distance de l'origine et du point $(1,1,1)$ au plan $x+2y+2z=9$.
\end{enumerate}
\ifcorr
\textbf{Corrigé et barème.} \textbf{1.} $r=2$ \bareme{2} ; $\cos\theta=\frac{\sqrt3}2$, $\sin\theta=-\frac12$, $\theta=-\frac\pi6$ \bareme{4}. \textbf{2.} $u\cdot v=4$ \bareme{3} ; $\|u\|=\sqrt6$, $\|v\|=\sqrt5$, $\cos\gamma=\frac4{\sqrt{30}}$ \bareme{3} ; positif, donc aigu \bareme{1}. \textbf{3.} $\sqrt{1+4+4}=3$ ; origine : $\frac93=3$ \bareme{3} ; $(1,1,1)$ : $\frac{|1+2+2-9|}3=\frac43$ \bareme{4}. Seuil conseillé : $14/20$.
\fi

\subsection{Grille d'auto-évaluation}
$\square$ je détermine le quadrant avant de donner $\theta$ ; $\square$ je sais quel angle est $\varphi$ dans mon cours ; $\square$ j'écris le plan sous la forme $ax+by+cz=d$ avant la distance ; $\square$ je vérifie $u\ne0$ avant de projeter ; $\square$ je n'oublie pas $r$ ou $\rho^2\sin\varphi$ dans les volumes.
\clearpage
\section{Calcul matriciel}
\renewcommand{\titrecourant}{Semaine 6}

\begin{boussole}
\textbf{Objectif de la semaine.} Multiplier des matrices en sachant que l'ordre compte, calculer un déterminant et dire ce qu'il mesure, décider si une matrice est inversible, calculer son inverse, et résoudre un système linéaire en distinguant les trois situations : solution unique, infinité, aucune.\\
\textbf{À la fin, je sais :}
\begin{enumerate}
 \item vérifier la compatibilité des tailles puis calculer $AB$ ligne par colonne ;
 \item calculer $\det$ en $2\times2$ et en $3\times3$ et l'interpréter (aire, volume, coplanarité) ;
 \item calculer $A^{-1}$ ($2\times2$ par la formule, $3\times3$ par comatrice) et le vérifier par $AA^{-1}=I$ ;
 \item résoudre $AX=b$ quand $\det A\ne0$ et analyser le cas $\det A=0$.
\end{enumerate}
\textbf{Temps estimé} : 5 à 6 heures. L'annexe sur les vecteurs propres est un prolongement \horsprog\ (la diagonalisation n'est pas au programme du module) ; elle n'est pas évaluée.
\end{boussole}

\subsection{Test d'entrée (5 minutes)}
\begin{enumerate}
 \item Calculer $2\cdot5-3\cdot4$.
 \item Résoudre $\begin{cases}x+y=3\\x-y=1\end{cases}$.
 \item Quelle est l'aire du parallélogramme de côtés $(2,0)$ et $(1,3)$ ?
 \item Que vaut $\frac1a\cdot a$ ? Pour quel $a$ le nombre $\frac1a$ n'existe-t-il pas ?
 \item Produit scalaire de $(1,2)$ et $(3,4)$.
\end{enumerate}
\textit{Réponses : $-2$ ; $x=2$, $y=1$ ; $6$ ; $1$, $a=0$ ; $11$.}

\subsection{Cours}
\begin{definition}[Produit matriciel]
Si $A$ est de taille $m\times p$ et $B$ de taille $p\times n$, alors $AB$ est de taille $m\times n$ avec $(AB)_{ij}=\sum_{k=1}^pA_{ik}B_{kj}$. Le produit n'existe que si le nombre de colonnes de $A$ égale le nombre de lignes de $B$.
\end{definition}
\begin{propriete}[Règles]
$(AB)C=A(BC)$ ; $A(B+C)=AB+AC$ ; $AI=IA=A$ ; en général $AB\ne BA$ ; et $AB=0$ n'implique pas $A=0$ ou $B=0$.
\end{propriete}
\begin{definition}[Déterminant]
$\det\begin{pmatrix}a&b\\c&d\end{pmatrix}=ad-bc$. En $3\times3$, développement selon une ligne avec le signe en damier $\begin{smallmatrix}+&-&+\\-&+&-\\+&-&+\end{smallmatrix}$ : $\det A=a_{11}C_{11}+a_{12}C_{12}+a_{13}C_{13}$, où $C_{ij}=(-1)^{i+j}\times$ (déterminant $2\times2$ obtenu en retirant ligne $i$ et colonne $j$).
\end{definition}
\begin{propriete}[Interprétation et produit]
En $2\times2$, $|\det|$ est l'aire du parallélogramme construit sur les colonnes ; en $3\times3$, le volume du parallélépipède. Le déterminant est nul si et seulement si les colonnes sont linéairement dépendantes (coplanaires en dimension 3). De plus $\det(AB)=\det A\det B$.
\end{propriete}
\begin{propriete}[Inverse]
$A$ carrée est inversible si et seulement si $\det A\ne0$. En $2\times2$ : $A^{-1}=\frac1{ad-bc}\begin{pmatrix}d&-b\\-c&a\end{pmatrix}$. En général $A^{-1}=\frac1{\det A}\,{}^t\mathrm{Com}(A)$ où $\mathrm{Com}(A)$ est la matrice des cofacteurs. Et $(AB)^{-1}=B^{-1}A^{-1}$ (ordre inversé).
\end{propriete}
\begin{demonstration}[formule $2\times2$]
$\begin{pmatrix}a&b\\c&d\end{pmatrix}\begin{pmatrix}d&-b\\-c&a\end{pmatrix}=\begin{pmatrix}ad-bc&0\\0&ad-bc\end{pmatrix}=(ad-bc)I$ ; on divise par $ad-bc\ne0$. De même dans l'autre sens.
\end{demonstration}
\begin{propriete}[Système carré]
Soit $AX=b$ avec $A$ carrée. Si $\det A\ne0$ : solution unique $X=A^{-1}b$. Si $\det A=0$ : soit aucune solution, soit une infinité.
\end{propriete}
\begin{methode}
\textbf{Résoudre un système.} (1) Écrire $AX=b$. (2) Calculer $\det A$. (3) Si $\ne0$ : $X=A^{-1}b$, puis \emph{vérifier} en reportant. (4) Si $=0$ : élimination (combinaisons de lignes) pour trancher entre incompatible et infinité de solutions.
\end{methode}
\begin{piege}
(i) Multiplier terme à terme. (ii) Écrire $(AB)^{-1}=A^{-1}B^{-1}$. (iii) Oublier un signe dans le damier. (iv) Appliquer la formule de l'inverse sans tester $\det A\ne0$. (v) Conclure \og une seule solution \fg{} sans avoir calculé $\det A$.
\end{piege}

\subsection{Exemples travaillés et cas limites}
\begin{exemple}[non-commutativité]
$A=\begin{pmatrix}1&2\\0&1\end{pmatrix}$, $B=\begin{pmatrix}1&0\\3&1\end{pmatrix}$.\\
\emph{Solution.} $AB=\begin{pmatrix}7&2\\3&1\end{pmatrix}$, $BA=\begin{pmatrix}1&2\\3&7\end{pmatrix}$. Ces deux matrices sont différentes, bien que les deux produits existent.
\end{exemple}
\begin{cas}[produit nul, tailles incompatibles]
$\begin{pmatrix}1&0\\0&0\end{pmatrix}\begin{pmatrix}0&0\\0&1\end{pmatrix}=\begin{pmatrix}0&0\\0&0\end{pmatrix}$ alors qu'aucun facteur n'est nul : on ne peut pas \og simplifier \fg{} par une matrice. Par ailleurs, une matrice $2\times3$ multipliée par une matrice $2\times3$ n'a pas de produit (3 colonnes contre 2 lignes) ; mais $2\times3$ par $3\times2$ donne une matrice $2\times2$.
\end{cas}
\begin{exemple}[déterminant $2\times2$ et $3\times3$]
$\det\begin{pmatrix}3&1\\5&2\end{pmatrix}=6-5=1$. Pour $M=\begin{pmatrix}1&2&0\\0&1&3\\2&0&1\end{pmatrix}$, selon la première ligne : $1\cdot(1\cdot1-3\cdot0)-2\cdot(0\cdot1-3\cdot2)+0=1+12=13$.
\end{exemple}
\begin{cas}[déterminant nul]
$N=\begin{pmatrix}2&0&1\\1&3&2\\1&1&1\end{pmatrix}$ : $\det N=2(3-2)-0+1(1-3)=2-2=0$. Les colonnes sont coplanaires : $C_1=-C_2+2C_3$ (vérification : $-(0,3,1)+2(1,2,1)=(2,1,1)$). Le volume du parallélépipède est nul et $N$ n'est pas inversible.
\end{cas}
\begin{exemple}[inverse $2\times2$]
$\begin{pmatrix}2&1\\5&3\end{pmatrix}^{-1}=\frac1{6-5}\begin{pmatrix}3&-1\\-5&2\end{pmatrix}=\begin{pmatrix}3&-1\\-5&2\end{pmatrix}$. Contrôle : $\begin{pmatrix}2&1\\5&3\end{pmatrix}\begin{pmatrix}3&-1\\-5&2\end{pmatrix}=\begin{pmatrix}1&0\\0&1\end{pmatrix}$. $\checkmark$
\end{exemple}
\begin{exemple}[inverse $3\times3$]
Pour $M$ ci-dessus ($\det M=13$), le calcul des cofacteurs donne $13\,M^{-1}=\begin{pmatrix}1&-2&6\\6&1&-3\\-2&4&1\end{pmatrix}$. Contrôle sur la première ligne de $M$ : $(1,2,0)\cdot(1,6,-2)=13$ (colonne 1 de la matrice ci-dessus), et $(1,2,0)\cdot(-2,1,4)=0$. $\checkmark$
\end{exemple}
\begin{cas}[matrice singulière]
$\begin{pmatrix}2&4\\1&2\end{pmatrix}$ : $\det=4-4=0$. La seconde ligne est la moitié de la première. La formule de l'inverse divise par $0$ : pas d'inverse. Elle écrase le plan sur une droite (aire nulle).
\end{cas}
\begin{exemple}[système à solution unique]
$M\begin{pmatrix}x\\y\\z\end{pmatrix}=\begin{pmatrix}3\\4\\3\end{pmatrix}$. Comme $\det M=13\ne0$, $X=\frac1{13}\begin{pmatrix}1&-2&6\\6&1&-3\\-2&4&1\end{pmatrix}\begin{pmatrix}3\\4\\3\end{pmatrix}=\frac1{13}\begin{pmatrix}13\\13\\13\end{pmatrix}=\begin{pmatrix}1\\1\\1\end{pmatrix}$. Contrôle : $M(1,1,1)^T=(3,4,3)^T$. $\checkmark$
\end{exemple}
\begin{cas}[trois situations pour $\det A=0$]
Les systèmes $\begin{cases}x+y=2\\2x+2y=4\end{cases}$ et $\begin{cases}x+y=2\\2x+2y=5\end{cases}$ ont la même matrice de déterminant $0$. Le premier a une infinité de solutions ($y=2-x$) ; le second n'en a aucune (la seconde équation donnerait $x+y=\frac52$). Seul le second membre les distingue. Et $\begin{cases}x+y=2\\x-y=0\end{cases}$ ($\det=-2\ne0$) n'a que $(1,1)$.
\end{cas}

\subsection{Exercices}
\begin{exo}[A]\label{c6e1}
Pour $C=\begin{pmatrix}2&1\\0&3\end{pmatrix}$ et $D=\begin{pmatrix}1&4\\2&0\end{pmatrix}$, calculer $CD$ et $DC$.
\end{exo}
\begin{exo}[A]\label{c6e2}
Calculer le déterminant puis l'inverse de $\begin{pmatrix}3&1\\5&2\end{pmatrix}$.
\end{exo}
\begin{exo}[A]\label{c6e3}
Calculer $\det\begin{pmatrix}1&2&3\\0&1&4\\5&6&0\end{pmatrix}$.
\end{exo}
\begin{exo}[B]\label{c6e4}
Pour quelles valeurs de $m$ la matrice $\begin{pmatrix}1&m\\m&4\end{pmatrix}$ est-elle inversible ? Donner l'inverse en fonction de $m$.
\end{exo}
\begin{exo}[B]\label{c6e5}
Résoudre $\begin{pmatrix}2&1\\5&3\end{pmatrix}X=\begin{pmatrix}1\\2\end{pmatrix}$.
\end{exo}
\begin{exo}[B]\label{c6e6}
Un parallélogramme est construit sur $u=(3,1)$ et $v=(5,2)$. Calculer son aire. Que se passe-t-il pour $v=(6,2)$ ?
\end{exo}
\begin{exo}[C]\label{c6e7}
Discuter les solutions de $\begin{cases}x+y=2\\2x+2y=k\end{cases}$ selon $k$.
\end{exo}
\begin{exo}[C]\label{c6e8}
Avec $A$ et $B$ de l'exemple de non-commutativité, vérifier $\det(AB)=\det A\det B$ et $(AB)^{-1}=B^{-1}A^{-1}$.
\end{exo}
\begin{exo}[C]\label{c6e9}
Montrer que si $A$ est inversible et $AB=0$, alors $B=0$.
\end{exo}

\ifcorr
\subsection{Corrigés détaillés}
\begin{corrige}{\ref{c6e1}}
$CD=\begin{pmatrix}2\cdot1+1\cdot2&2\cdot4+0\\0+3\cdot2&0\end{pmatrix}=\begin{pmatrix}4&8\\6&0\end{pmatrix}$ ; $DC=\begin{pmatrix}2&1+12\\4&2\end{pmatrix}=\begin{pmatrix}2&13\\4&2\end{pmatrix}$. Donc $CD\ne DC$.
\end{corrige}
\begin{corrige}{\ref{c6e2}}
$\det=6-5=1$ ; $A^{-1}=\begin{pmatrix}2&-1\\-5&3\end{pmatrix}$. Contrôle : $\begin{pmatrix}3&1\\5&2\end{pmatrix}\begin{pmatrix}2&-1\\-5&3\end{pmatrix}=\begin{pmatrix}1&0\\0&1\end{pmatrix}$. $\checkmark$
\end{corrige}
\begin{corrige}{\ref{c6e3}}
Selon la première ligne : $1\cdot(0-24)-2\cdot(0-20)+3\cdot(0-5)=-24+40-15=1$.
\end{corrige}
\begin{corrige}{\ref{c6e4}}
$\det=4-m^2=(2-m)(2+m)$ : inversible si et seulement si $m\ne\pm2$. Alors $A^{-1}=\frac1{4-m^2}\begin{pmatrix}4&-m\\-m&1\end{pmatrix}$. Contrôle $m=0$ : $\frac14\begin{pmatrix}4&0\\0&1\end{pmatrix}=\mathrm{diag}(1,\frac14)$, inverse de $\mathrm{diag}(1,4)$. $\checkmark$ Pour $m=2$, les lignes $(1,2)$ et $(2,4)$ sont proportionnelles.
\end{corrige}
\begin{corrige}{\ref{c6e5}}
$\det=1$, inverse $\begin{pmatrix}3&-1\\-5&2\end{pmatrix}$ : $X=\begin{pmatrix}3-2\\-5+4\end{pmatrix}=\begin{pmatrix}1\\-1\end{pmatrix}$. Contrôle : $2-1=1$, $5-3=2$. $\checkmark$
\end{corrige}
\begin{corrige}{\ref{c6e6}}
Aire $=|\det\begin{pmatrix}3&5\\1&2\end{pmatrix}|=|6-5|=1$. Avec $v=(6,2)=2u$ : $\det=6-6=0$, les vecteurs sont colinéaires, le parallélogramme est aplati, d'aire nulle.
\end{corrige}
\begin{corrige}{\ref{c6e7}}
$\det=0$ pour tout $k$. La seconde équation s'écrit $x+y=\frac k2$ : si $k=4$, elle coïncide avec la première, infinité de solutions $(x,2-x)$ ; si $k\ne4$, incompatible avec $x+y=2$, aucune solution.
\end{corrige}
\begin{corrige}{\ref{c6e8}}
$AB=\begin{pmatrix}7&2\\3&1\end{pmatrix}$ : $\det=7-6=1=\det A\det B=1\cdot1$. $B^{-1}=\begin{pmatrix}1&0\\-3&1\end{pmatrix}$, $A^{-1}=\begin{pmatrix}1&-2\\0&1\end{pmatrix}$, $B^{-1}A^{-1}=\begin{pmatrix}1&-2\\-3&7\end{pmatrix}$, qui est l'inverse de $AB$ (formule $2\times2$ : $\begin{pmatrix}1&-2\\-3&7\end{pmatrix}$). $\checkmark$ Le produit $A^{-1}B^{-1}=\begin{pmatrix}7&-2\\-3&1\end{pmatrix}$ est différent.
\end{corrige}
\begin{corrige}{\ref{c6e9}}
Multiplions $AB=0$ à gauche par $A^{-1}$ : $A^{-1}AB=IB=B=A^{-1}\cdot0=0$.
\end{corrige}
\else\stepcounter{subsection}\fi

\subsection{Contrôle de fin de semaine (30 minutes, 20 points)}
\begin{enumerate}
 \item \bareme{6} Calculer $PQ$ et $QP$ pour $P=\begin{pmatrix}2&1\\0&3\end{pmatrix}$ et $Q=\begin{pmatrix}1&4\\2&0\end{pmatrix}$.
 \item \bareme{7} Montrer que $R=\begin{pmatrix}2&1\\5&3\end{pmatrix}$ est inversible et calculer $R^{-1}$.
 \item \bareme{7} Résoudre $RX=\begin{pmatrix}1\\2\end{pmatrix}$ et vérifier.
\end{enumerate}
\ifcorr
\textbf{Corrigé et barème.} \textbf{1.} $PQ=\begin{pmatrix}4&8\\6&0\end{pmatrix}$ \bareme{3}, $QP=\begin{pmatrix}2&13\\4&2\end{pmatrix}$ \bareme{3}. \textbf{2.} $\det R=1\ne0$ \bareme{3} ; $R^{-1}=\begin{pmatrix}3&-1\\-5&2\end{pmatrix}$ \bareme{4}. \textbf{3.} $X=R^{-1}b=(1,-1)^T$ \bareme{4} ; vérification $2-1=1$, $5-3=2$ \bareme{3}. Seuil conseillé : $14/20$.
\fi

\subsection*{Annexe \horsprog\ : valeurs et vecteurs propres}
Un vecteur $v\ne0$ est propre pour la valeur $\lambda$ si $Av=\lambda v$. Une valeur $\lambda$ est propre si et seulement si $\det(A-\lambda I)=0$ : cette condition assure l'existence d'un vecteur propre non nul, sans déterminer ce vecteur. Pour $S=\begin{pmatrix}2&1\\1&2\end{pmatrix}$ : $\det(S-\lambda I)=(2-\lambda)^2-1=(\lambda-1)(\lambda-3)$, donc $\lambda=1$ ou $\lambda=3$ ; vecteurs propres $(1,-1)$ pour $1$ et $(1,1)$ pour $3$. Contrôle : $S(1,1)^T=(3,3)^T$. Cas limite : $\begin{pmatrix}0&-1\\1&0\end{pmatrix}$ (rotation d'un quart de tour) n'a aucune valeur propre réelle ($\lambda^2+1=0$).

\subsection{Grille d'auto-évaluation}
$\square$ je vérifie les tailles avant de multiplier ; $\square$ je n'écris jamais $AB=BA$ sans preuve ; $\square$ je calcule $\det A$ avant de parler d'inverse ou d'unicité ; $\square$ je contrôle $AA^{-1}=I$ ; $\square$ pour $\det A=0$, je compare les seconds membres avant de conclure.
\clearpage
\section{Fonctions de deux variables, intégrales multiples, équations différentielles}
\renewcommand{\titrecourant}{Semaine 7}

\begin{boussole}
\textbf{Objectif de la semaine.} Dériver une fonction de deux variables par rapport à chacune, en déduire un plan tangent et une propagation d'incertitude ; calculer une intégrale double (ordre, changement de variables) ; résoudre une équation différentielle linéaire d'ordre 1 et d'ordre 2 à coefficients constants.\\
\textbf{À la fin, je sais :}
\begin{enumerate}
 \item calculer $f_x$, $f_y$ et écrire le plan tangent ;
 \item estimer une erreur sur $f(x,y)$ à partir d'erreurs sur $x$ et $y$ ;
 \item calculer $\iint_D f$ en choisissant le bon ordre ou le bon système de coordonnées (jacobien) ;
 \item résoudre $y'+ay=b$ et $y''+by'+cy=0$ avec conditions initiales, en distinguant les trois cas du discriminant.
\end{enumerate}
\textbf{Temps estimé} : 6 heures.
\end{boussole}

\subsection{Test d'entrée (5 minutes)}
\begin{enumerate}
 \item Dériver $x^2\e^{3x}$ par rapport à $x$.
 \item Calculer $\int_0^1\int_0^2 xy\,\dd y\,\dd x$ en intégrant d'abord en $y$.
 \item Racines de $r^2+3r+2=0$ ?
 \item Primitive de $\e^{-2t}$ ?
 \item Écrire $x^2+y^2$ en coordonnées polaires.
\end{enumerate}
\textit{Réponses : $(2x+3x^2)\e^{3x}$ ; $1$ ; $-1$ et $-2$ ; $-\frac12\e^{-2t}$ ; $r^2$.}

\subsection{Cours}
\subsubsection*{Dérivées partielles}
\begin{definition}[Dérivées partielles]
$f_x(a,b)=\lim_{h\to0}\frac{f(a+h,b)-f(a,b)}h$ : on dérive en $x$ en gelant $y$ ; de même $f_y$. Le plan tangent à la surface $z=f(x,y)$ au point $(a,b)$ est $z=f(a,b)+f_x(a,b)(x-a)+f_y(a,b)(y-b)$, lorsque $f$ est différentiable en ce point.
\end{definition}
\begin{propriete}[Schwarz et incertitudes]
Si $f_{xy}$ et $f_{yx}$ existent et sont continues, elles sont égales. Pour de petites variations : $\Delta f\approx f_x\Delta x+f_y\Delta y$ ; l'estimation au premier ordre de l'erreur absolue est $|f_x||\Delta x|+|f_y||\Delta y|$. Ce n'est pas une borne rigoureuse de l'erreur exacte sans contrôle du reste.
\end{propriete}
\subsubsection*{Intégrales doubles}
\begin{propriete}[Fubini]
Si $f$ est continue sur un domaine $D=\{a\le x\le b,\ g_1(x)\le y\le g_2(x)\}$, alors $\iint_Df=\int_a^b\Big(\int_{g_1(x)}^{g_2(x)}f(x,y)\,\dd y\Big)\dd x$. On peut échanger l'ordre en réécrivant les bornes de $D$ \emph{avec l'autre variable extérieure}.
\end{propriete}
\begin{propriete}[Changement de variables]
Si $(x,y)=\Phi(u,v)$ est un difféomorphisme $C^1$ entre les intérieurs de $D'$ et $D$ (jacobien non nul), à des frontières de mesure nulle près, alors $\iint_Df\,\dd x\,\dd y=\iint_{D'}f(\Phi(u,v))\,|J|\,\dd u\,\dd v$ où $J=\det\frac{\partial(x,y)}{\partial(u,v)}$. En polaires, $J=r$ ; pour $(x,y)=(u+v,u-v)$, $J=-2$.
\end{propriete}
\subsubsection*{Équations différentielles}
\begin{propriete}[Ordre 1]
Les solutions de $y'+ay=b$ ($a\ne0$) sont $y(t)=\frac ba+C\e^{-at}$, avec $C=y(0)-\frac ba$. Si $a=0$ : $y(t)=y(0)+bt$.
\end{propriete}
\begin{propriete}[Ordre 2 homogène]
Pour $y''+by'+cy=0$, soit $r^2+br+c=0$ de discriminant $\Delta=b^2-4c$. Si $\Delta>0$ (racines $r_1,r_2$) : $y=A\e^{r_1t}+B\e^{r_2t}$. Si $\Delta=0$ (racine double $r$) : $y=(A+Bt)\e^{rt}$. Si $\Delta<0$ ($r=\alpha\pm i\omega$) : $y=\e^{\alpha t}(A\cos\omega t+B\sin\omega t)$. Les deux constantes sont fixées par $y(0)$ et $y'(0)$.
\end{propriete}
\begin{demonstration}[cas $\Delta>0$ : ces fonctions sont solutions]
Si $r_i^2+br_i+c=0$, alors $(\e^{r_it})''+b(\e^{r_it})'+c\,\e^{r_it}=(r_i^2+br_i+c)\e^{r_it}=0$ ; l'équation est linéaire, donc toute combinaison $A\e^{r_1t}+B\e^{r_2t}$ est solution. Le fait que ce soient \emph{toutes} les solutions est admis dans ce cahier.
\end{demonstration}
\begin{piege}
(i) Dériver $f$ en $x$ sans geler $y$. (ii) Écrire les bornes de l'ordre échangé sans redessiner $D$. (iii) Oublier le jacobien $r$. (iv) Écrire $\e^{rt}$ pour une racine double sans le facteur $(A+Bt)$. (v) Croire que l'existence de $f_x,f_y$ garantit la continuité.
\end{piege}

\subsection{Exemples travaillés et cas limites}
\begin{exemple}[dérivées partielles]
$f(x,y)=x^2y+\e^{xy}$.\\
\emph{Solution.} $f_x=2xy+y\e^{xy}$ ; $f_y=x^2+x\e^{xy}$. Contrôle de Schwarz : $f_{xy}=2x+\e^{xy}+xy\e^{xy}=f_{yx}$. $\checkmark$
\end{exemple}
\begin{cas}[existence des dérivées partielles sans régularité]
(a) $f=\sqrt{x^2+y^2}$ en $(0,0)$ : $f(x,0)=|x|$ n'est pas dérivable en $0$, donc $f_x(0,0)$ n'existe pas (cône à pointe, sans plan tangent). (b) $f(x,y)=\frac{xy}{x^2+y^2}$, $f(0,0)=0$ : $f(x,0)=0$ et $f(0,y)=0$, donc $f_x(0,0)=f_y(0,0)=0$. Pourtant $f(x,x)=\frac12$ pour $x\ne0$ : $f$ n'est même pas continue en $0$. Les dérivées partielles ne suffisent pas à garantir un plan tangent.
\end{cas}
\begin{exemple}[plan tangent]
$f=x^2+y^2$ en $(1,2)$ : $f=5$, $f_x=2$, $f_y=4$. Plan : $z=5+2(x-1)+4(y-2)$. Estimation de $f(1{,}1;1{,}9)$ : $5+0{,}2-0{,}4=4{,}8$ ; valeur exacte $1{,}21+3{,}61=4{,}82$.
\end{exemple}
\begin{exemple}[propagation d'incertitude]
Le volume d'un cylindre $V=\pi r^2h$ est connu avec $r$ à $2\,\%$ et $h$ à $1\,\%$. Estimation au premier ordre de l'erreur relative maximale : $2\cdot2\,\%+1\,\%=5\,\%$. Ce n'est pas une borne exacte : pour des variations positives maximales, $(1{,}02)^2(1{,}01)-1=0{,}050804$, soit $5{,}0804\,\%$.
\end{exemple}
\begin{cas}[estimation au premier ordre et erreur réelle]
Pour $Q=\frac rh$, avec les mêmes incertitudes, l'estimation au premier ordre est $2\,\%+1\,\%=3\,\%$ : on additionne les valeurs absolues des contributions. La hausse exacte maximale est $1{,}02/0{,}99-1\approx3{,}0303\,\%$. L'erreur \emph{effective} peut être plus petite si les erreurs se compensent (par exemple $r$ trop grand et $h$ trop grand en même temps) ; l'estimation maximale au premier ordre ne donne pas l'erreur effective.
\end{cas}
\begin{exemple}[intégrale double sur un triangle]
$D=\{0\le y\le x\le1\}$ ; calculer $\iint_D(x+y)$.\\
\emph{Solution.} $\int_0^1\int_0^x(x+y)\,\dd y\,\dd x=\int_0^1\big(x^2+\frac{x^2}2\big)\dd x=\int_0^1\frac{3x^2}2\dd x=\frac12$.
\end{exemple}
\begin{cas}[l'ordre d'intégration décide de la faisabilité]
Le même domaine $D=\{0\le y\le x\le1\}$ donne deux écritures. Dans l'ordre $\dd y$ puis $\dd x$ : $\int_0^1\int_0^x\e^{x^2}\,\dd y\,\dd x=\int_0^1x\e^{x^2}\,\dd x=\frac{\e-1}2$, calcul immédiat. Dans l'ordre $\dd x$ puis $\dd y$ : $\int_0^1\int_y^1\e^{x^2}\,\dd x\,\dd y$, où $\e^{x^2}$ n'a pas de primitive élémentaire : le calcul est bloqué. Les deux écritures valent la même chose ; il faut choisir celle qu'on sait intégrer.
\end{cas}
\begin{exemple}[coordonnées polaires]
$\iint_{x^2+y^2\le R^2}(x^2+y^2)\,\dd x\,\dd y=\int_0^{2\pi}\int_0^Rr^2\cdot r\,\dd r\,\dd\theta=\frac{\pi R^4}2$. Et $\iint_{\R^2}\e^{-(x^2+y^2)}=2\pi\int_0^\infty r\e^{-r^2}\dd r=\pi$.
\end{exemple}
\begin{cas}[jacobien nul]
Pour $(x,y)=(u+v,u-v)$, $|J|=2$ : le carré unité du plan $(u,v)$ devient un losange d'aire $2$. Pour $(x,y)=(u,2u)$ (non bijective sur un domaine d'aire non nulle), $J=0$ : l'image est une droite, d'aire nulle ; en polaires $r=0$ est le seul point où $J$ s'annule, de mesure nulle.
\end{cas}
\begin{exemple}[ordre 1]
$y'+2y=6$, $y(0)=0$ : solution particulière $3$ ; $y=3+C\e^{-2t}$, $C=-3$ : $y(t)=3-3\e^{-2t}$.
\end{exemple}
\begin{cas}[$a=0$ et blow-up]
$y'=6$ ($a=0$) : $y=y(0)+6t$, la formule $\frac ba$ n'existe pas. L'équation non linéaire $y'=y^2$, $y(0)=1$, a pour solution $y=\frac1{1-t}$, qui tend vers l'infini en $t=1$ : la solution n'existe pas pour tout $t$ (explosion en temps fini).
\end{cas}
\begin{exemple}[ordre 2, trois cas]
Avec $y(0)=1$, $y'(0)=0$. (a) $y''+3y'+2y=0$ : $r=-1,-2$, $y=2\e^{-t}-\e^{-2t}$. (b) $y''+2y'+y=0$ : $r=-1$ double, $y=(1+t)\e^{-t}$. (c) $y''+2y'+5y=0$ : $r=-1\pm2i$, $y=\e^{-t}\big(\cos2t+\frac12\sin2t\big)$. Contrôle de (a) : $y(0)=2-1=1$, $y'(0)=-2+2=0$. $\checkmark$
\end{exemple}
\begin{cas}[sans amortissement]
$y''+4y=0$ : $r=\pm2i$, $\alpha=0$, $y=\cos2t$ avec les mêmes conditions : oscillation pure, sans décroissance. $y''-4y=0$ : $r=\pm2$ et $y=\cosh2t=\frac{\e^{2t}+\e^{-2t}}2$, qui croît sans borne.
\end{cas}

\subsection{Exercices}
\begin{exo}[A]\label{c7e1}
Calculer $f_x$ et $f_y$ pour $f(x,y)=x^3y-\e^{y}$, puis leurs valeurs en $(1,0)$.
\end{exo}
\begin{exo}[A]\label{c7e2}
Écrire le plan tangent à $z=xy$ au point $(2,3)$ et estimer $f(2{,}1;2{,}9)$.
\end{exo}
\begin{exo}[A]\label{c7e3}
Calculer $\iint_{[0,1]\times[0,2]}xy\,\dd x\,\dd y$.
\end{exo}
\begin{exo}[B]\label{c7e4}
Calculer $\iint_Dxy$ sur le triangle $D=\{x\ge0,\ y\ge0,\ x+y\le1\}$.
\end{exo}
\begin{exo}[B]\label{c7e5}
Calculer $\int_0^1\int_y^1\sin(x^2)\,\dd x\,\dd y$ en échangeant l'ordre.
\end{exo}
\begin{exo}[B]\label{c7e6}
Résoudre $y'+y=2$, $y(0)=5$.
\end{exo}
\begin{exo}[C]\label{c7e7}
Un café est à $90^\circ$C dans une pièce à $20^\circ$C. Après $10$ minutes il est à $60^\circ$C. Avec la loi $T'=-k(T-20)$, déterminer $k$ puis le temps pour atteindre $30^\circ$C.
\end{exo}
\begin{exo}[C]\label{c7e8}
Résoudre $y''+4y'+4y=0$, $y(0)=0$, $y'(0)=1$.
\end{exo}
\begin{exo}[C]\label{c7e9}
Déduire de $\iint_{\R^2}\e^{-(x^2+y^2)}=\pi$ la valeur de $\int_{-\infty}^{+\infty}\e^{-x^2}\,\dd x$.
\end{exo}

\ifcorr
\subsection{Corrigés détaillés}
\begin{corrige}{\ref{c7e1}}
$f_x=3x^2y$, $f_y=x^3-\e^y$. En $(1,0)$ : $f_x=0$, $f_y=1-1=0$ : le point $(1,0)$ est critique.
\end{corrige}
\begin{corrige}{\ref{c7e2}}
$f(2,3)=6$, $f_x=y=3$, $f_y=x=2$ : $z=6+3(x-2)+2(y-3)$. En $(2{,}1;2{,}9)$ : $6+0{,}3-0{,}2=6{,}1$ ; exact : $6{,}09$.
\end{corrige}
\begin{corrige}{\ref{c7e3}}
Les bornes sont constantes : $\int_0^1x\,\dd x\cdot\int_0^2y\,\dd y=\frac12\cdot2=1$.
\end{corrige}
\begin{corrige}{\ref{c7e4}}
$\int_0^1x\int_0^{1-x}y\,\dd y\,\dd x=\int_0^1\frac{x(1-x)^2}2\dd x=\frac12\int_0^1(x-2x^2+x^3)\dd x=\frac12\big(\frac12-\frac23+\frac14\big)=\frac1{24}$.
\end{corrige}
\begin{corrige}{\ref{c7e5}}
Le domaine est $0\le y\le x\le1$. En échangeant : $\int_0^1\int_0^x\sin(x^2)\,\dd y\,\dd x=\int_0^1x\sin(x^2)\,\dd x=\big[-\frac{\cos(x^2)}2\big]_0^1=\frac{1-\cos1}2\approx0{,}2298$.
\end{corrige}
\begin{corrige}{\ref{c7e6}}
Solution constante $2$ ; $y=2+C\e^{-t}$, $C=5-2=3$ : $y=2+3\e^{-t}$. Vérification : $y'=-3\e^{-t}$ et $y'+y=2$. $\checkmark$
\end{corrige}
\begin{corrige}{\ref{c7e7}}
$T=20+70\e^{-kt}$. À $t=10$ : $70\e^{-10k}=40$, $\e^{-10k}=\frac47$, $k=\frac{\ln(7/4)}{10}\approx0{,}0560\ \mathrm{min}^{-1}$. $T=30$ : $\e^{-kt}=\frac17$, $t=\frac{\ln7}k=\frac{10\ln7}{\ln(7/4)}\approx34{,}8$ minutes.
\end{corrige}
\begin{corrige}{\ref{c7e8}}
$r^2+4r+4=(r+2)^2$ : racine double $-2$, $y=(A+Bt)\e^{-2t}$. $y(0)=A=0$ ; $y'(0)=B-2A=B=1$. Donc $y=t\e^{-2t}$. Vérification : $y''+4y'+4y=(-4+4t)+4(1-2t)+4t=0$. $\checkmark$
\end{corrige}
\begin{corrige}{\ref{c7e9}}
Par Fubini sur le produit : $\iint\e^{-x^2}\e^{-y^2}=\big(\int\e^{-x^2}\big)^2=\pi$, et l'intégrale est positive : $\int_{-\infty}^{+\infty}\e^{-x^2}\,\dd x=\sqrt\pi$.
\end{corrige}
\else\stepcounter{subsection}\fi

\subsection{Contrôle de fin de semaine (30 minutes, 20 points)}
\begin{enumerate}
 \item \bareme{6} Pour $f(x,y)=x^2\e^y$, calculer $f_x$, $f_y$ et vérifier $f_{xy}=f_{yx}$.
 \item \bareme{7} Calculer $\iint_{0\le y\le x\le1}2xy\,\dd y\,\dd x$.
 \item \bareme{7} Résoudre $y''+2y'+y=0$, $y(0)=1$, $y'(0)=0$.
\end{enumerate}
\ifcorr
\textbf{Corrigé et barème.} \textbf{1.} $f_x=2x\e^y$, $f_y=x^2\e^y$ \bareme{4} ; $f_{xy}=2x\e^y=f_{yx}$ \bareme{2}. \textbf{2.} Intérieure : $\int_0^x2xy\,\dd y=x^3$ \bareme{4} ; $\int_0^1x^3=\frac14$ \bareme{3}. \textbf{3.} Racine double $-1$ \bareme{2} ; $y=(A+Bt)\e^{-t}$, $A=1$ \bareme{2} ; $y'(0)=B-A=0$, $B=1$, $y=(1+t)\e^{-t}$ \bareme{3}. Seuil conseillé : $14/20$.
\fi

\subsection{Grille d'auto-évaluation}
$\square$ je gèle l'autre variable en dérivant ; $\square$ je dessine le domaine avant d'écrire les bornes ; $\square$ je pense au jacobien en polaires ; $\square$ j'écris l'équation caractéristique et son discriminant ; $\square$ je vérifie la solution en la reportant dans l'équation.

\end{document}
